\documentclass[11pt,oneside,openany]{book}
\usepackage[T1]{fontenc}
\usepackage{lmodern}
\usepackage{amsmath,amssymb,amsthm,mathtools}
\usepackage[letterpaper,margin=1.08in]{geometry}
\usepackage{microtype,enumitem,booktabs,xcolor,needspace}
\usepackage[unicode,colorlinks=true,linkcolor=blue!52!black,citecolor=teal!60!black,urlcolor=blue!52!black]{hyperref}
\usepackage{bookmark,cleveref}
\newtheorem{theorem}{Theorem}[chapter]
\newtheorem{lemma}[theorem]{Lemma}
\newtheorem{proposition}[theorem]{Proposition}
\newtheorem{corollary}[theorem]{Corollary}
\theoremstyle{definition}
\newtheorem{definition}[theorem]{Definition}
\newtheorem{example}[theorem]{Example}
\theoremstyle{remark}
\newtheorem{remark}[theorem]{Remark}
\newcommand{\R}{\mathbb R}
\newcommand{\HH}{\mathbb H}
\newcommand{\dd}{\mathrm d}
\DeclareMathOperator{\supp}{supp}
\DeclareMathOperator{\id}{id}
\allowdisplaybreaks[1]
\setlength{\emergencystretch}{2em}
\setcounter{tocdepth}{1}
\hypersetup{pdftitle={Smooth Manifolds and Intrinsic Calculus},
 pdfauthor={MathGenome},pdfsubject={P1-04 Unified Mathematical Notes}}

\begin{document}
\frontmatter
\hypersetup{pageanchor=false}
\begin{titlepage}
\thispdfpagelabel{Cover}
\centering
\vspace*{0.18\textheight}
{\large\scshape MathGenome\par}
\vspace{2em}
{\Huge\bfseries Smooth Manifolds and\par Intrinsic Calculus\par}
\vspace{1.6em}
{\Large Unified mathematical notes\par}
\vfill
{\large P1-04\par}
\vspace{0.5em}
{7 November 2026\par}
\end{titlepage}
\hypersetup{pageanchor=true}

\chapter*{Scope and conventions}
\addcontentsline{toc}{chapter}{Scope and conventions}

These notes develop the calculus of abstract smooth manifolds from charts. Local normal forms and global localization lead to Sard's theorem, Whitney embedding and approximation, and transversality. Vector bundles, flows, differential forms, and Frobenius give the tools for integration. The final chapters establish Stokes's theorem, basic de Rham cohomology and degree, and construct the integration map on smooth singular cochains. The full comparison with continuous singular cohomology is stated at the comparison interface. Riemannian, Lie group, and symplectic constructions appear where they illustrate the underlying calculus.

Manifolds are finite-dimensional, Hausdorff, and second countable, with a fixed dimension unless indicated otherwise. Smooth means \(C^\infty\) over \(\R\). For \(n\geq1\) put \(\HH^n=\{x\in\R^n:x_n\geq0\}\). A boundary chart is relatively open in this model. Smoothness at a face is defined using local extensions to Euclidean open sets. The zero-dimensional model is a point and has empty boundary. We distinguish a manifold's intrinsic boundary, the frontier of a subset in an ambient space, and the local faces in a corner chart. The support of a section is the closure in its base of its nonzero locus.

The differential of a map at \(p\) is \(\dd f_p\), also denoted \(T_pf\). The bracket of vector fields is their commutator on functions. A field \(X\) has local flow \(\varphi_t\) on its domain of definition. Lie differentiation is the derivative at zero of \(\varphi_t^*\). Interior multiplication \(\iota_X\) inserts \(X\) in the first slot. The evaluation of a wedge of one-forms is the determinant of their values on the ordered vectors. Boundary orientations are outward-normal first. For cochains the differential raises degree, and oriented simplices have the alternating face boundary. Coefficients of differential forms and de Rham groups are real.

The prerequisites are elementary finite-dimensional linear and multilinear algebra, basic general topology, Euclidean differential and integral calculus, Lebesgue measure and Fubini's theorem, and local ODE existence, uniqueness and smooth dependence in Euclidean open sets. The Euclidean inverse, implicit and constant-rank theorems, Euclidean change of variables, and the regular Lindel\"of paracompactness theorem are used in their usual stated settings. The relevant MathGenome volumes are \cite{d05,d10,d04,d11,d13}. The two main textbook sources are \cite{lee,tu}; substantive corrections follow the authors' errata \cite{leeerrata,tuerrata}. The result map accompanying these notes records the primary statement for each of the 74 writing items and the dependencies that discharge it.

{\small\tableofcontents}
\mainmatter
\chapter{Smooth structures, boundary, and localization}

\section{Topological models}

An \(n\)-dimensional topological manifold is a Hausdorff, second-countable
space that has charts into relatively open subsets of \(\HH^n\), for \(n>0\).
The model \(\R^n\) can be used when all points are interior. Dimension zero
uses the one-point model. A chart is a homeomorphism onto its image, with
open domain in the manifold. Throughout, the word \emph{boundary} refers to
the intrinsic boundary unless an ambient frontier is specified.

\begin{proposition}[Basic topology of a manifold]\label{p104:M01.01}
A topological manifold is locally compact and locally path connected.
It admits a countable chart cover, is paracompact, and has a compact
exhaustion \(K_1\subset\operatorname{int}K_2\subset\cdots\), where
\(M=\bigcup_{j\geq1}\operatorname{int}K_j\). Its connected components
are open and path connected, and there are at most countably many of them.
\end{proposition}

\begin{proof}
In a model chart, intersect sufficiently small open balls with the model.
These relative balls are convex and have compact relative closures
inside the chart. Such neighborhoods can be chosen inside any given
neighborhood of the point. Their transported closures are compact and
closed in \(M\). This gives local compactness, local path connectedness,
and the shrinking property used in the regularity axiom. Second
countability implies Lindel\"ofness, so regular Lindel\"of
paracompactness applies \cite{d10}. Taking a countable subcover from all
precompact chart neighborhoods gives a sequence \((V_j)\) with compact
closures. For the exhaustion, start with a compact set containing
\(\overline V_1\) in its interior. Inductively cover the compact set
\(K_j\cup\overline V_{j+1}\) by finitely many precompact chart
neighborhoods and take the union of their closures for \(K_{j+1}\).
A finite chart cover can be repeated to use the same construction.
When \(M=\varnothing\), take every \(K_j\) empty. Finally, path
components are open because each point has small path-connected
neighborhoods. They are also closed, since their complements are
unions of path components, and so are exactly the connected components.
Every nonempty component contains a member of a fixed countable basis,
and disjoint components contain distinct members.
\end{proof}

\section{Atlases and the intrinsic boundary}

On a relatively open subset of a Euclidean half-space, a map into
Euclidean space is smooth if at each point it agrees locally with a
smooth map on a Euclidean open neighborhood. The same convention
applies componentwise to a map with values in a half-space, with the
extension allowed to take values outside it. In positive dimension the
relative open domain is the closure of its interior locally. Hence two
local extensions agreeing on the domain have equal derivatives there.
Their derivatives agree on its Euclidean interior and then on its
boundary by continuity. This works for derivatives of every available
order, and for \(C^r\) extensions of finite order \(r\).

\begin{definition}[Smooth structures]\label{p104:M01.02}
An atlas is a chart cover. It is smooth if on every overlap the chart
transition and its inverse are smooth on their model domains. A smooth
structure is a maximal smooth atlas, and a smooth manifold is a
topological manifold equipped with such a structure.
\end{definition}

Each smooth atlas extends uniquely to a maximal one. Indeed, adjoin all
charts compatible with every original chart. Two adjoined charts have
smooth transition on their overlap because at each point of the overlap
their transition factors locally through an original chart. Any smooth
atlas containing the original one is contained in this enlarged atlas.
The same argument gives a construction from parametrizations of a set.
If the coordinate overlap domains are relatively open and the
transitions are homeomorphisms, declare \(U\) open when its intersection
with every chart domain has relatively open coordinates. Then every
chart is a homeomorphism and has open domain. Check Hausdorffness
separately, and use a countable chart subcover for second countability
when it exists. Smooth compatibility then gives the structure. See
\cite[Chapter~1]{lee} and \cite[Sections~5--6]{tu}.

\begin{proposition}[Products and standard constructions]\label{p104:M01.03}
Open subsets carry the restricted smooth structure. The product of two
manifolds without boundary has the product structure, and a product
with exactly one boundary factor carries the half-space product
structure. In particular, spheres, tori, and real projective spaces
have the following natural smooth structures.
\end{proposition}

\begin{proof}
For an open subset, restrict chart domains. In a product, the
coordinate changes are products of the original transitions. In the
case with one boundary factor, permute the coordinates so that its
boundary coordinate is last. Finite products and open subspaces
preserve the two separation and countability hypotheses. A product of
two boundary factors has a corner model at a pair of boundary points.
The product charts there will be used in
\cref{p104:M10.06}.

For the sphere \(S^n\subset\R^{n+1}\), omit a coordinate on each of its
positive or negative hemispheres. The inverse inserts
\(\pm\sqrt{1-\|x\|^2}\), which is smooth in the open unit ball.
Coordinate changes are obtained by taking components of these
parametrizations. The torus is \((S^1)^k\) with the product atlas.
The map \(\R^k/\mathbb Z^k\to(S^1)^k\) defined by the periodic angle
coordinates is a homeomorphism, as the quotient is compact and the
target is Hausdorff. On sufficiently small coordinate boxes the angle
map has a smooth inverse, giving the corresponding quotient charts.

Write \(\R P^n\) as the antipodal quotient of \(S^n\). Distinct
orbits can be separated by disjoint small metric neighborhoods of the
two finite orbits. These neighborhoods are antipodally invariant
because the metric is. The quotient map is
open, since the saturation of an open subset is the union of it and
its antipodal image, so a countable basis descends. On the chart where
the \(i\)-th coordinate is nonzero, use the ratios \(x_j/x_i\) for
\(j\neq i\). The inverse inserts \(1\) in position \(i\) and takes
the line of the resulting vector. Continuity follows also by
normalizing this representative to the unit sphere. Transitions are
rational expressions with a nonvanishing denominator. These
descriptions include \(S^0\) and \(\R P^0\).
\end{proof}

\begin{theorem}[Invariance and structure of the boundary]\label{p104:M01.04}
For \(n\geq1\), the condition that a point have last coordinate zero in
a half-space chart is independent of the smooth chart. The resulting
subset \(\partial M\) is closed in \(M\) and is naturally an
\((n-1)\)-dimensional smooth manifold without boundary. The complement
is the manifold interior. For \(n=0\) set \(\partial M=\varnothing\).
\end{theorem}

\begin{proof}
Let \(f\) and \(g\) be inverse chart transitions and choose their local
Euclidean extensions near corresponding points. Shrink the ambient
neighborhoods to form both compositions. On the relative chart
domains the compositions are identities. Differentiating on their
Euclidean interiors and taking limits proves that the extension
derivatives are inverse linear maps at the given points. Suppose a
point of the boundary face is sent to an interior point. Near the
interior image, \(g\) is a map on a Euclidean open set. Its derivative
is invertible, so the Euclidean inverse function theorem says its
image contains a Euclidean neighborhood of the original point. This
is impossible because its image is contained in the half-space and
the original point is on the face. Interchanging \(f\) and \(g\)
proves the other implication. The same argument applies to a
transition into a chart modeled on an open Euclidean set.

The interior is open in every chart, hence \(\partial M\) is closed.
The charts restricted to the face have relatively open domains in
\(\partial M\) and images open in the face \(\R^{n-1}\).
Their transitions are restrictions of the extensions to the
face, with values in the face, and are smooth. Thus these
restricted charts form the indicated smooth atlas. They also
show that the inclusion \(\partial M\to M\) is smooth.
\end{proof}

This argument uses smooth chart compatibility. It establishes the
smooth boundary invariance directly, without presupposing the
corresponding topological invariance result
\cite[Theorem~1.46]{lee}. See also \cite[Section~22]{tu}.

\section{Maps and extension conventions}

\begin{proposition}[Local tests for smooth maps]\label{p104:M01.05}
A continuous map \(f:M\to N\) is smooth if each point of \(M\) has a
source chart whose image under \(f\) lies in a target chart, with
smooth coordinate representation there. The test is independent of
the charts, respects restriction and gluing on open subsets, and is
closed under composition. These assertions hold for \(C^r\) maps of
\(C^r\) atlases as well.
\end{proposition}

\begin{proof}
Let \(x\) lie in two source chart domains. Continuity allows shrinking
the source domain around \(x\) so that its image is in the overlap of
the two target domains being compared. The new coordinate formula is
the old one composed on the two sides with the appropriate
transitions. At boundary points choose their extensions and shrink
once more so that the composition is defined in Euclidean
neighborhoods. This proves independence of charts and the
composition assertion. The same local definition proves restriction
and gluing. The latter includes continuity because continuity is
local on an open cover.
\end{proof}

A diffeomorphism is a smooth bijection with smooth inverse.
For example, \(x\mapsto x^3\) is a smooth bijection of \(\R\) with
inverse not differentiable at \(0\). A diffeomorphism transports
the atlas, and preserves the intrinsic boundary by the proof of
\cref{p104:M01.04}. For a subset \(A\subset M\), we call
a map \(A\to\R^k\) smooth on \(A\) when it extends to a smooth
map near each point of \(A\), agreeing on the intersection with
\(A\). An analogous local extension convention applies to a
boundaryless manifold target, but the global extension argument
below is for vector values. See \cite[Chapter~2, especially p.~45]{lee},
\cite[p.~45]{leeerrata}, and \cite[Section~6]{tu}.

\section{Bump functions and smooth partitions of unity}

We use the following local construction. The function
\(\theta(t)=e^{-1/t}\) for \(t>0\) and \(0\) for \(t\leq0\)
is smooth, since each one-sided derivative is a polynomial in
\(1/t\) times \(e^{-1/t}\), with limit zero at the origin.
The quotient
\[
  \eta(t)=\frac{\theta(1-t)}{\theta(1-t)+\theta(t)}
\]
is smooth, equals \(1\) on \(t\leq0\), and equals \(0\) on
\(t\geq1\). Given two concentric balls of radii \(0<a<b\), apply
\(\eta\) to \((\|x-x_0\|^2-a^2)/(b^2-a^2)\).
This gives a nonnegative smooth function that is \(1\) on the
smaller closed ball and supported in the larger closed ball.
Restriction to the half-space gives the same relative construction
at a boundary point. If the larger closed relative ball lies in
a chart domain, the function transported to \(M\) extends
smoothly by zero outside that domain.

\begin{theorem}[Smooth localization]\label{p104:M01.06}
Every open cover of a smooth manifold admits a smooth partition of
unity subordinate to it. In addition, if \(A\) is closed and
\(A\subset U\) with \(U\) open, there is a smooth function
\(0\leq\chi\leq1\) equal to \(1\) on a neighborhood of \(A\), with
\(\supp\chi\subset U\). If \(A\) is compact, \(\chi\) can be chosen
with compact support.
\end{theorem}

\begin{proof}
Let \(K_j\) be the exhaustion in \cref{p104:M01.01}, and put
\(K_0=K_{-1}=\varnothing\). The compact shells
\(A_j=K_j\setminus\operatorname{int}K_{j-1}\) are contained in
the open bands
\(B_j=\operatorname{int}K_{j+1}\setminus K_{j-2}\).
For each point of \(A_j\), choose two small coordinate balls as
above, with the larger closed ball in \(B_j\) and inside one
member of the given cover. Finitely many of the smaller open
balls cover \(A_j\). The transported bump functions \(b_i\)
are locally finite, because a neighborhood contained in
\(\operatorname{int}K_m\) is disjoint from \(B_j\) for
\(j\geq m+2\). The shells cover \(M\), so
\(b=\sum_i b_i\) is smooth and strictly positive on \(M\).
The functions \(b_i/b\) sum to \(1\) and have the required
closed supports. They may be grouped according to the original
cover index. The union of each corresponding locally finite
family of closed supports is closed and contained in its
assigned cover set. The empty manifold has the empty partition.

For the cutoff use the cover \(\{U,M\setminus A\}\), grouping the
weights assigned to \(U\). The other weights have a closed union
of supports disjoint from \(A\), so the resulting sum equals \(1\)
on an open neighborhood of \(A\). Its support is contained in
the closed union of supports assigned to \(U\). For compact \(A\)
first choose an open neighborhood of \(A\) with compact closure
inside \(U\) by a finite union of small coordinate balls, and
replace \(U\) by that smaller neighborhood in the cutoff argument.
\end{proof}

The support condition in this construction is stronger than pointwise
vanishing outside the relevant open set. It permits extensions by zero
and the common refinements needed for integration. See
\cite[Theorem~2.23 and Proposition~2.25]{lee} and
\cite[Section~13 and Appendix~C]{tu}.

\begin{theorem}[Extension, exhaustion, and zero sets]\label{p104:M01.07}
Let \(M\) be a smooth manifold, \(A\subset M\) closed, and
\(A\subset U\) open.
\begin{enumerate}[label=\textup{(\alph*)}]
\item A smooth map \(f:A\to\R^k\) in the local extension sense has
  a smooth extension to \(M\) supported in \(U\).
\item There is a smooth nonnegative \(\rho:M\to\R\) whose sublevel
  sets \(\{\rho\leq c\}\) are compact for every real \(c\).
\item There is a smooth nonnegative \(g:M\to\R\) with
  \(g^{-1}(0)=A\).
\end{enumerate}
\end{theorem}

\begin{proof}
For (a), take local extensions on open neighborhoods in \(U\)
of the points of \(A\), and use them with \(M\setminus A\) as an
open cover. Weight the extensions with a subordinate partition,
using the zero function on \(M\setminus A\). A weight times its
local extension is smooth after extension by zero, since the
weight is supported inside the extension domain. The weighted
sum is locally finite and agrees with \(f\) on \(A\), where only
the extensions agreeing with \(f\) contribute. Its support is
contained in the locally finite union of closed supports of the
weights assigned to sets in \(U\).

For (b), take cutoffs \(\chi_j\) equal to \(1\) near \(K_j\)
and supported in \(\operatorname{int}K_{j+1}\). Set
\(\rho=\sum_{j\geq1}(1-\chi_j)\). Near any point only finitely
many summands can be nonzero, since some
\(\operatorname{int}K_j\) contains a neighborhood of the point.
If \(x\notin K_{N+1}\), the first \(N\) terms are \(1\).
For \(N>c\) the sublevel set is therefore a closed subset of
\(K_{N+1}\).

For (c), first let \(C\) be closed in an open set
\(V\subset\R^n\). Cover \(V\setminus C\) by countably many balls
with compact closures inside \(V\setminus C\), and choose
nonnegative smooth Euclidean bumps \(h_j\) supported on those
closed balls and positive in their interiors. For example use
\(\theta(b_j^2-\|x-x_j\|^2)\) for the ball of radius \(b_j\)
centered at \(x_j\). Choose constants
\[
 a_j=\frac{2^{-j}}{1+
   \max_{|\alpha|\leq j}\|\partial^\alpha h_j\|_\infty}.
\]
For each fixed derivative order the tail of
\(\sum_j a_j\partial^\alpha h_j\) converges uniformly, and the
finitely many initial terms cause no difficulty. The sum
\(h=\sum_j a_jh_j\) is smooth with those derivatives. To see
the differentiation claim, apply the fundamental theorem of
calculus along short coordinate segments to each finite sum,
pass to uniform limits on the segment, and induct on the order.
Every summand vanishes on \(C\), while every point outside \(C\)
in \(V\) has a positive summand. Empty complements and
zero-dimensional charts have the same assertion directly.

In a manifold chart write the image as \(V\cap\HH^n\) with
\(V\) open in \(\R^n\), or just \(V\) for an interior chart.
The image of \(A\) in the chart is relatively closed in the chart
image, which is itself
closed in \(V\) in the half-space case. The Euclidean
construction thus gives an actual ambient smooth function,
so the restriction is smooth even at the boundary. Multiply
the resulting chart functions by weights of a partition
subordinate to these charts and extend the products by zero.
Their locally finite nonnegative sum vanishes on \(A\).
At a point off \(A\) at least one weight is positive and its
chart function is positive, proving the equality of zero sets.
\end{proof}

These constructions apply to arbitrary closed sets and to noncompact
manifolds \cite[Lemma~2.26, Proposition~2.28, and Theorem~2.29]{lee}.
\chapter{Tangent spaces, tangent maps, and the tangent bundle}

\section{Derivations and curve velocities}

Write \(\mathcal C_p^\infty(M)\) for the real algebra of germs
of smooth real-valued functions at \(p\). Two functions define
the same germ when they agree on a neighborhood of \(p\).
A derivation at \(p\) is a linear map
\(D:\mathcal C_p^\infty(M)\to\R\) satisfying
\[
 D(fg)=f(p)Dg+g(p)Df.
\]
It kills constants by applying this identity to \(1\cdot1\).
Let \(\mathfrak m_p\) be the ideal of germs vanishing at \(p\).
Derivations are naturally the linear dual of
\(\mathfrak m_p/\mathfrak m_p^2\). Restriction to
\(\mathfrak m_p\) kills products of vanishing germs, and a
linear functional on this quotient becomes a derivation by
applying it to the class of \(f-f(p)\).

\begin{theorem}[Models of the tangent space]\label{p104:M02.01}
For every point of a smooth \(n\)-manifold, the derivations at
\(p\) form an \(n\)-dimensional vector space \(T_pM\).
Coordinates \(x=(x^1,\ldots,x^n)\) give the basis
\(\partial/\partial x^i|_p\). The coordinates of \(D\) are
\(D(x^i)\), and a curve \(c\) with \(c(0)=p\) determines the
derivation \(f\mapsto (f\circ c)'(0)\). Every derivation is
represented by a curve, allowing the parameter interval to
have \(0\) as either endpoint at a boundary point. Equality
of curve velocities is independent of the chart.
\end{theorem}

\begin{proof}
For a germ \(f\), choose coordinates centered at \(a=x(p)\),
an ambient smooth extension of the coordinate representative,
and a small Euclidean ball about \(a\). The segment formula on
this ball yields the identity of germs
\[
  f-f(p)=\sum_i(x^i-a^i)g_i,\qquad
  g_i(x)=\int_0^1
    \partial_i(f\circ x^{-1})(a+t(x-a))\,\dd t ,
\]
where the integrand uses the extension when needed. The
functions \(g_i\) are smooth, by differentiating the smooth
integrand on the compact parameter interval. Their values at
\(p\) are the coordinate partial derivatives, which are
extension independent by Chapter~1. The derivation rule
therefore gives
\[
  Df=\sum_i D(x^i)\,\partial_i(f\circ x^{-1})(a).
\]
Conversely the right side defines a derivation for any list of
real coefficients. Testing on the coordinate functions shows
uniqueness. Thus \(T_pM\) has the stated basis. The corresponding
coordinate formula on overlaps is the usual derivative of
the transition, by the Euclidean chain rule. A curve with
coordinate derivative \(v\) acts on germs by that same formula,
so the equality of velocities is intrinsic. At an interior
point a straight coordinate curve realizes any \(v\).
At a boundary point, the straight ray in one of the two
parameter directions realizes \(v\), since its normal
coordinate is \(tv_n\) and a one-sided interval can be chosen
with \(tv_n\geq0\). When \(v_n=0\) a two-sided segment is
available. For \(n=0\), all germs are constant and \(T_pM=0\).
\end{proof}

See \cite[pp.~50--55 and 68--73]{lee},
\cite[Section~8]{tu}, and \cite[Problem~3-8]{leeerrata} for the
boundary curve convention. One can also use derivations on
globally defined functions. Such a derivation kills any \(f\)
which vanishes near \(p\). To see this, multiply \(f\) by a
cutoff equal to \(1\) near \(p\) and supported where \(f=0\), and apply the
derivation rule to the zero product. It therefore depends
only on germs. Conversely a germ extends to a global smooth
function after multiplication by a cutoff, so the two
definitions are equivalent.

\begin{proposition}[The boundary tangent subspace and cones]\label{p104:M02.02}
For \(p\in\partial M\), the differential of the inclusion
identifies \(T_p\partial M\) with the hyperplane \(v_n=0\)
in a half-space chart. Forward velocities of curves staying
in \(M\) are exactly \(v_n\geq0\). Backward velocities are
exactly \(v_n\leq0\), and two-sided velocities are exactly
\(T_p\partial M\). In particular \(T_pM\) has dimension \(n\)
at a boundary point as well.
\end{proposition}

\begin{proof}
On a boundary chart the inclusion of the face is
\(u\mapsto(u,0)\). Its derivative has image \(v_n=0\).
The same formula follows from the action on germs. A
curve defined for \(t\geq0\) satisfies \(x_n(c(t))\geq0\)
and \(x_n(c(0))=0\), hence its right derivative is
nonnegative. The left derivative is nonpositive for a
curve defined on \(t\leq0\). The straight curves in the
preceding proof realize every permitted vector.

For independence of the half-space chart, write a smooth
coordinate transition as \(F\). Its last component is zero
on the boundary hyperplane, by \cref{p104:M01.04}.
Its tangential derivatives of that component are zero,
and its inward normal derivative is nonnegative because
it maps the half-space into the half-space. That normal
derivative cannot be zero, since the full derivative is
invertible. Thus the last coordinate of the transformed
vector is a strictly positive multiple of the original
last coordinate. This proves independence of the cones
and identifies the normal quotient line
\(T_pM/T_p\partial M\) with a line having a well-defined
inward and outward half.
\end{proof}

In particular, using only two-sided curves through a boundary
point would discard the normal direction. See
\cite[Proposition~3.23]{lee} and \cite[Section~22.4]{tu}.

\section{Differentiating maps}

\begin{proposition}[Differentials and the chain rule]\label{p104:M02.03}
For smooth \(f:M\to N\), the formula
\[
  \bigl(\dd f_p(v)\bigr)(h)=v(h\circ f)
  \qquad(h\in\mathcal C_{f(p)}^\infty(N))
\]
defines a linear map \(T_pM\to T_{f(p)}N\).
In coordinates it is the Jacobian of the local coordinate
representative. It respects identity maps and satisfies
\[
  \dd(g\circ f)_p=\dd g_{f(p)}\circ\dd f_p .
\]
For the products of \cref{p104:M01.03}, the tangent space is the direct sum of the two
tangent spaces and the differential of a pair of maps is
the pair of differentials. In particular, a diffeomorphism
induces fiberwise linear isomorphisms. The coordinate
definition and chain rule also apply to \(C^r\) maps of
\(C^r\) manifolds, \(r\geq1\).
\end{proposition}

\begin{proof}
The composite \(h\circ f\) is well-defined as a germ at \(p\),
and its value at \(p\) is \(h(f(p))\). These facts give the
derivation rule on the target algebra. Evaluating on target
coordinate functions and using \cref{p104:M02.01} proves the
Jacobian description. It also shows that if \(v\) is realized
by \(c\), then \(\dd f_p(v)\) is realized by \(f\circ c\).
The chain rule in germs follows from associativity of
composition. For products, the coordinates of the two
projections are exactly the two blocks of coordinates,
so their differentials identify \(T_{(p,q)}(M\times N)\)
with \(T_pM\oplus T_qN\). The statement on diffeomorphisms
follows by applying the chain rule to both inverse
compositions.

At finite regularity define the tangent space from
coordinate vector representatives using the transition
derivatives. These derivatives are independent of the
ambient extensions at boundary points, and their
cocycle identity follows from the Euclidean \(C^1\)
chain rule and continuity on the face. Hence the
coordinate construction and chain rule need only
\(C^1\) regularity. The transition maps for the total
tangent space and the induced tangent map have
regularity \(C^{r-1}\). This finite-regularity argument
does not require the algebraic germ model.
\end{proof}

See \cite[pp.~55--65]{lee} and \cite[Sections~8 and 10]{tu}.

\section{From individual tangent spaces to a bundle}

\begin{theorem}[The tangent bundle and the tangent functor]\label{p104:M02.04}
The set
\[
  TM=\coprod_{p\in M} T_pM
\]
has a canonical smooth manifold structure. Its projection
\(\pi_M:TM\to M\) is locally a product with the vector
space \(\R^n\). The fiber coordinate transitions are linear
isomorphisms depending smoothly on the base point.
The zero section and the fiberwise vector operations
are smooth. Each smooth map \(f:M\to N\) induces a
smooth map \(Tf:TM\to TN\) covering \(f\), with
\(T\id_M=\id_{TM}\) and \(T(g\circ f)=Tg\circ Tf\).
For a manifold with boundary,
\(\partial(TM)=\pi_M^{-1}(\partial M)\).
\end{theorem}

\begin{proof}
Above a chart \(x:U\to V\), use the bijection
\[
  v\in T_pM \longmapsto \bigl(x(p),\,\dd x_p(v)\bigr)
  \in V\times\R^n.
\]
Above overlaps the transition is
\[
  (u,v)\longmapsto
  \bigl(G(u),\,\dd G_u(v)\bigr),\qquad
  G=y\circ x^{-1},
\]
with inverse given by the same formula for \(G^{-1}\).
Thus it is smooth, including at the face by local
extension. These charts determine the topology as in
\cref{p104:M01.02}. Points over different base points
have disjoint neighborhoods by Hausdorffness of the
base and continuity of the projection. Two points in
one fiber can be separated in the same product chart.
A countable chart cover of \(M\) gives a countable
cover of \(TM\) by second-countable open charts,
so \(TM\) is second countable. The displayed charts
also give the boundary identity and the usual smooth
coordinate descriptions of the zero section, addition
on pairs in the same fiber, and scalar multiplication.

The map \(Tf\) sends \(v\in T_pM\) to \(\dd f_p(v)\).
In source and target charts it is
\((u,v)\mapsto(F(u),\dd F_u(v))\), where \(F\) is the
coordinate representative of \(f\). This is smooth
and covers \(f\), and its functoriality follows from
\cref{p104:M02.03}. For dimension zero the fibers
are singletons, and the construction reduces to the
discrete base.
\end{proof}

Thus \(TM\) is a smooth vector bundle in the sense
formalized in \cref{p104:M05.01}. See
\cite[pp.~65--73]{lee} and \cite[Sections~10 and 12]{tu}.

\section{Comparison with an embedded model}

\begin{proposition}[Embedded tangent spaces]\label{p104:M02.05}
Suppose a smooth manifold \(M\) is identified with a subset
of \(\R^N\) and its local parametrizations into \(\R^N\)
are smooth maps of rank \(n=\dim M\). Write \(j\) for the
inclusion. For a parametrization \(\psi\) sending \(u\)
to \(p\), the image
\[
   \dd j_p(T_pM)=\dd\psi_u(\R^n)
\]
is independent of \(\psi\), and \(\dd j_p\) identifies
\(T_pM\) linearly with this image. At a boundary point,
the image of \(T_p\partial M\) is the image of the
tangent hyperplane of the parametrizing half-space.
\end{proposition}

\begin{proof}
In the chart on \(M\) underlying \(\psi\), the map \(j\)
has coordinate representative \(\psi\). Thus the asserted
image is exactly the Jacobian image by
\cref{p104:M02.03}, and the rank hypothesis makes the
map onto that image an isomorphism. On overlaps the
derivative of the reparametrization is invertible, so
composing by it leaves the image unchanged. The boundary
statement follows by applying \cref{p104:M02.02} and
then the same chain rule to the inclusion of the face.
The full image includes the normal direction within \(M\).
Let \(\gamma\) be a smooth two-sided ambient curve with
\(\gamma(0)=j(p)\) and a continuous local lift through \(p\).
Such a lift is automatic when \(M\) has the subspace topology
and the curve stays in \(j(M)\). Choose a coordinate projection
\(P:\R^N\to\R^n\) for which \(P\dd\psi_u\) is invertible.
For a smooth Euclidean extension \(\widetilde\psi\) near \(u\),
the inverse function theorem gives a smooth local inverse of
\(P\widetilde\psi\). By continuity, the lift stays in the
chart where this inverse applies for parameters near \(0\).
Its coordinate expression there is
\((P\widetilde\psi)^{-1}\circ P\circ\gamma\). The lift is
therefore smooth, so at a boundary point \cref{p104:M02.02}
and the chain rule put the ambient
velocity in \(\dd j_p(T_p\partial M)\).
\end{proof}

This is the tangent-space comparison used for the
embedded calculus in \cite[Chapter~14]{d04}. The argument
also explains why comparing an intrinsic and an embedded
integration theorem requires matching their boundary and
tangent data in addition to comparing the underlying sets.
\chapter{Local normal forms, submanifolds, and smooth descent}

\section{Rank and normal forms}

The rank of a \(C^1\) map at a point is the rank of its
differential there. A map is an immersion where this
differential is injective, and a submersion where it is
surjective. These properties are invariant under chart
changes by \cref{p104:M02.03}. The following statements
separate open models from the boundary situations in which
an extension can be used.

\begin{theorem}[Inverse and constant-rank normal forms]\label{p104:M03.01}
Let \(f:M^m\to N^n\) be \(C^r\), \(r\geq1\).
\begin{enumerate}[label=\textup{(\alph*)}]
\item At interior points \(p\) and \(f(p)\), if \(\dd f_p\)
  is invertible, then \(f\) is a local \(C^r\) diffeomorphism.
  If the rank is constantly \(k\) near \(p\) and both
  charts are interior, local coordinates give
  \((u,v)\mapsto(u,0)\) with \(u\in\R^k\).
  In particular, immersions and submersions have
  the local inclusion and projection forms.
\item If \(f\) is an immersion at a boundary point of \(M\)
  and \(f(p)\) is an interior point of \(N\), source and
  target coordinates give the same inclusion form,
  with the source restricted to its half-space.
\item If \(p\) and \(f(p)\) are boundary points of equal
  dimensional manifolds, \(f\) carries nearby boundary
  points to the boundary, and \(\dd f_p\) is invertible,
  then \(f\) is a local \(C^r\) diffeomorphism of the
  half-space models.
\end{enumerate}
\end{theorem}

\begin{proof}
For (a), first apply the Euclidean inverse function
theorem in charts in the invertible case. For constant
rank, a linear change of coordinates puts a nonsingular
\(k\)-minor in the first \(k\) variables and components.
For a representative \(F\), change source coordinates
by the local inverse of
\((x,y)\mapsto(F_1(x,y),\ldots,F_k(x,y),y)\).
Then \(F\) takes the form \((u,G(u,v))\). Its derivative
has block matrix
\[
  \begin{pmatrix} I_k&0\\ D_uG&D_vG\end{pmatrix}.
\]
Since its rank is \(k\), \(D_vG=0\). On a small product
box the fundamental theorem along segments in \(v\)
implies \(G(u,v)=G(u,v_0)\). Subtract \(G(u,v_0)\) in the
last target coordinates. The cases \(k=0\), \(k=m\), and
\(k=n\) use the same argument with the empty blocks
omitted. For maximal rank at a point, the relevant minor
remains nonsingular locally, so the rank condition follows.

For (b) extend the coordinate representative \(F\) to a
Euclidean neighborhood. Pick a linear injection
\(A:\R^{n-m}\to\R^n\) complementary to
\(\operatorname{im}\dd F_p\). The map
\((u,w)\mapsto F(u)+Aw\) has invertible derivative at
\((p,0)\). Its local inverse gives the target
coordinates and makes the original map \(u\mapsto(u,0)\)
with its original half-space condition on \(u\).

For (c), in half-space coordinates the last component
of an extension \(\widetilde F\) vanishes for \(x_n=0\)
near \(p\). Factoring it by the integral form of the
fundamental theorem gives
\(\widetilde F_n(x',s)=s\,h(x',s)\), where \(h\) is
continuous and \(h(p)=\partial_n\widetilde F_n(p)\).
Tangential partial derivatives in this row vanish at
\(p\). Invertibility therefore implies \(h(p)\ne0\),
and preservation of \(s\geq0\) implies \(h(p)>0\).
Shrink so that \(h>0\). The Euclidean local
diffeomorphism furnished by the inverse function
theorem carries the nonnegative half of its source
exactly to the nonnegative half of its target
neighborhood, proving (c).
\end{proof}

See \cite[Theorems~4.5, 4.12, and 4.15]{lee},
\cite[Section~11 and Appendix~B]{tu}, and the Euclidean
results in \cite[Chapter~4]{d04}. The inclusion
\(\HH^n\hookrightarrow\R^n\) has invertible derivative
everywhere and is not a local diffeomorphism at a
boundary point into \(\R^n\). It explains why the
boundary hypotheses of the local inverse statement
must be explicit.

\section{Embeddings, sections, and quotient maps}

An embedding is an injective immersion that is a homeomorphism
onto its image equipped with the subspace topology. A map
between manifolds is proper if inverse images of compact sets
are compact.

\begin{proposition}[Embedding criteria]\label{p104:M03.02}
An immersion is locally an embedding, including when
either manifold has boundary. A proper injective immersion is an
embedding. In particular an injective immersion from a compact
manifold is an embedding.
\end{proposition}

\begin{proof}
If the target point is on a boundary, view the local
coordinate representative as a map into the ambient
Euclidean space of its half-space chart. The inclusion
forms of \cref{p104:M03.01} then show local injectivity
and the homeomorphism property onto the image in that
ambient space. The image has the same subspace topology
in the target model. To prove the global criterion, first
let \(A\) be a closed set in the source of a proper map.
If \(y\) is in the closure of \(f(A)\), first countability
gives \(a_j\in A\) with \(f(a_j)\to y\). The sequence
together with its limit is compact. Its inverse image
is compact by properness and contains all \(a_j\).
It is also first countable, so a subsequence converges
to a point \(a\in A\), and continuity gives \(f(a)=y\).
Thus \(f\) is closed. For an injective map this means
the inverse on its image is continuous, as required.
Compact source implies properness because compact
subsets of the Hausdorff target are closed.
\end{proof}

\begin{example}
If \(\alpha\notin\mathbb Q\), the map
\(\R\to(\R/\mathbb Z)^2\) given by
\(t\mapsto([t],[\alpha t])\) is an injective immersion.
Its derivative is the nonzero vector \((1,\alpha)\)
in local quotient charts. Equality at two parameters
would force their difference and its product with
\(\alpha\) both to be integers. Nevertheless there are
positive integers \(q_j\to\infty\) with
\(\alpha q_j\) tending to \(0\) modulo \(\mathbb Z\).
Indeed the pigeonhole principle applied to the first
\(N+1\) multiples of \(\alpha\) modulo \(1\) gives
\(1\leq q\leq N\) at distance at most \(1/N\) from an
integer after multiplication by \(\alpha\). No fixed
finite set of positive \(q\) can achieve these distances
as \(N\) increases. Thus the image sequence approaches
the image of \(0\), while its parameters do not.
\end{example}

\begin{proposition}[Sections and smooth descent]\label{p104:M03.03}
A smooth submersion between boundaryless manifolds has a
smooth local section through every source point and is an
open map. A surjective smooth map \(q:M\to N\) having
smooth local sections around every point of \(N\) is a
quotient map. For such \(q\), any smooth \(h:M\to P\)
constant on the fibers of \(q\) descends to a unique
smooth \(\bar h:N\to P\). In particular, two surjective
smooth maps with local sections and the same fibers give
canonically diffeomorphic target manifolds.
\end{proposition}

\begin{proof}
The projection normal form gives a section through a
given point by fixing the unused coordinates. That form
is also open on small coordinate neighborhoods, and hence
the whole submersion is open. For the general claim, if
\(q^{-1}(B)\) is open and \(s:U\to M\) is a local section,
then \(B\cap U=s^{-1}(q^{-1}(B))\) is open in \(U\).
These neighborhoods cover \(N\). The set map
\(\bar h(y)=h(x)\) for \(q(x)=y\) is well-defined by the
fiber condition. On \(U\) it is \(h\circ s\), which
proves smoothness. Surjectivity proves uniqueness.
Apply this construction to the two surjections in
both directions. The two descended maps are inverses
because their composites are identities after
precomposition with the surjections.
\end{proof}

Compare \cite[pp.~85--91]{lee} and \cite[Sections~7 and 11]{tu}.
For a quotient built in a later construction it is enough
to verify actual smooth local sections. Surjectivity of
a boundary differential alone does not imply their
existence. For example, \(\HH^n\hookrightarrow\R^n\) has
surjective differential at a boundary point but is
not even locally onto an open target neighborhood there.

\section{Slices, regular levels, and boundary}

An immersed submanifold \(S\subset M\) is a smooth
manifold on the underlying set \(S\) whose inclusion is
an injective immersion. Its topology can be strictly
finer than the subspace topology, as in the irrational
line example. It is embedded when the inclusion is an
embedding. The intrinsic boundary of a manifold \(S\)
with boundary and the frontier of its image in \(M\)
are different notions.

\begin{theorem}[Slice and restriction criteria]\label{p104:M03.04}
In a boundaryless \(n\)-manifold, an embedded
\(k\)-dimensional submanifold is locally a coordinate
slice, with a \(k\)-dimensional half-slice at any point
of its own boundary. Conversely a subset with such
local slice charts carries a unique embedded structure.
A smooth map restricts smoothly to an immersed
submanifold of its source. A smooth map with image in
an immersed submanifold of its target is
smooth as a map into the submanifold when it is
continuous for the submanifold topology. Finally,
a smooth real-valued function on an embedded
submanifold
extends to a neighborhood of it, and to the entire
ambient manifold when the submanifold is closed.
\end{theorem}

\begin{proof}
Near a point of an embedded submanifold, put the
inclusion into the local normal form of
\cref{p104:M03.01}. Since the embedding is a
homeomorphism onto its image, an open source
neighborhood can be cut out by an open ambient
neighborhood. Shrink the ambient chart this way
to exclude any other part of the submanifold.
The image is now precisely a slice, or a half-slice
if the original source chart was a boundary chart.
Conversely take these slices for the charts of the
subset. On overlaps their transitions are the
restriction of an ambient transition followed
by the coordinate projection along the slice.
They and their inverses are smooth by local
extension, which proves existence. The same
coordinate description proves uniqueness and
that the inclusion is an embedding.

Restriction of the domain is composition with the
smooth inclusion. For restriction of the codomain,
an immersed submanifold is locally embedded in its
own topology by \cref{p104:M03.02}. Continuity into
that topology allows shrinking the source to land
in this one local embedded piece. If the target
chart is a half-space chart, straighten this local
immersion in the Euclidean space containing the
target model, by \cref{p104:M03.01}. The Euclidean
straightening is not required to preserve the
target half-space. In the resulting coordinates,
the induced map into the immersed submanifold is
the ambient coordinate representative followed
by projection onto its slice coordinates. This is
smooth by the local extension convention. For an
embedded submanifold the required continuity is
automatic.

In a slice chart a smooth real-valued function on
the submanifold extends locally by ignoring the
normal variables. The same Euclidean straightening
provides this extension if the ambient chart is
a half-space chart, by restricting the Euclidean
extension to the ambient model afterwards. At a
half-slice first extend the function across its own
face. The resulting local extensions agree on
the submanifold where their domains overlap. Take a
smooth partition of unity on the union of the
extension domains and patch there, or add the
complement of the submanifold to the cover with
zero local function when the submanifold is closed.
The same proof works for \(\R^k\)-valued functions
componentwise.
\end{proof}

See \cite[pp.~98--115, especially Theorems~5.8, 5.27, and
  5.29 and Lemma~5.34, and Theorems~5.51 and 5.53]{lee}
and \cite[Sections~9 and 11]{tu}.

\begin{theorem}[Regular levels and their tangent spaces]\label{p104:M03.05}
Suppose \(F:M^m\to N^n\) is smooth and both manifolds
are without boundary. If \(y\) is a regular value
of \(F\), meaning that \(\dd F_p\) is surjective for
every \(p\in F^{-1}(y)\), then \(F^{-1}(y)\) is an
embedded submanifold of codimension \(n\) whenever
it is nonempty. Its tangent space at \(p\) is
\(\ker\dd F_p\) under the differential of the inclusion.
The fiber is closed, and the empty fiber is allowed.
Conversely, an embedded boundaryless submanifold is
locally the regular zero set of a submersion into a
Euclidean space of dimension its codimension.
\end{theorem}

\begin{proof}
At a point of the fiber use target coordinates centered
at \(y\) and the submersion form. The fiber is the
coordinate slice with the first \(n\) coordinates zero.
Use \cref{p104:M03.04} for its embedded structure.
In these coordinates the derivative of \(F\) is a
projection, whose kernel is exactly the tangent
image of the slice. The fiber is closed because
\(\{y\}\) is closed. If \(n=0\), the target is
discrete and the fiber is open as well, consistent
with codimension zero. If \(m<n\), a regular fiber
is necessarily empty. Conversely the normal
coordinate projection in a slice chart is a
submersion with the submanifold as its zero set
in that chart.
\end{proof}

See \cite[pp.~104--108 and 115--120]{lee} and
\cite[Section~11.3]{tu}. The Euclidean tangent
identification also agrees with the regular-level
construction of \cite[Chapter~14]{d04}.

\begin{theorem}[Boundary regularity and defining functions]\label{p104:M03.06}
\begin{enumerate}[label=\textup{(\alph*)}]
\item Let \(F:M\to N\) be smooth, where \(M\) may have
  boundary and \(N\) has empty boundary. If \(y\) is
  regular for both \(F\) and \(F|_{\partial M}\), then
  \(S=F^{-1}(y)\)
  has the embedded manifold structure with
  \(\partial S=S\cap\partial M\). At a boundary
  point its full tangent space is \(\ker\dd F_p\)
  and the boundary tangent space is
  \(\ker\dd F_p\cap T_p\partial M\).
\item Let \(M\) be without boundary. A closed
  codimension-zero embedded submanifold \(D\) with
  boundary is a regular domain. Its intrinsic
  interior and boundary agree with its ambient
  interior and frontier, respectively. Such a
  \(D\) admits a smooth function \(f:M\to\R\)
  with \(D=\{f\leq0\}\), \(\partial D=f^{-1}(0)\),
  and \(0\) a regular value. If \(D\) is compact,
  \(f\) can be chosen to have compact sublevel sets.
  Conversely a regular sublevel set has this
  embedded structure.
\end{enumerate}
\end{theorem}

\begin{proof}
For (a) the interior part follows from
\cref{p104:M03.05}. At \(p\in S\cap\partial M\),
take source coordinates \((x',s)\) with \(s\geq0\)
and target coordinates centered at \(y\). The
boundary regularity hypothesis says the target
coordinate components have an invertible
\(n\)-minor using only coordinates among \(x'\).
Use those target components as the first \(n\)
new source coordinates, keep the remaining
tangential coordinates, and keep \(s\) unchanged.
An ambient extension of this coordinate change
has invertible derivative. Since its last
coordinate is \(s\), it restricts to a local
half-space diffeomorphism. In the new chart
the fiber is the half-slice given by the first
\(n\) coordinates zero. These half-slice and
interior-slice charts glue because their
transitions come from the ambient transitions
followed by projection. The boundary and
tangent formulas now follow directly. When
the boundary fiber is nonempty, its regularity
forces \(n\leq\dim M-1\), so the displayed
boundary dimension is meaningful. If \(n=0\)
the fiber is an open and closed portion of \(M\).

For (b), the slice description of
\cref{p104:M03.04} has full dimension at interior
points and gives an actual half-space at intrinsic
boundary points. This proves the assertions about
ambient interior and frontier because \(D\) is
closed. At each boundary point use the negative
of the half-space coordinate as a local defining
function. Add the function \(-1\) on the interior
of \(D\) and \(+1\) on \(M\setminus D\) to the
family. A subordinate smooth partition of unity
gives a global weighted sum. All terms contributing
at an interior point of \(D\) have negative sign, while
all at an exterior point have positive sign.
At a boundary point all participating local
defining functions vanish. Their differentials
are nonzero covectors annihilating the same
boundary tangent space and negative on the same
inward half of the quotient normal line.
Thus the derivative of the sum at that point
is a positive-weight sum of covectors with
the same strict sign on an inward vector. It
is nonzero. Derivatives of the weights contribute
zero there. If \(D\) is compact, take a cutoff
\(\chi\) equal to \(1\) near \(D\) with compact
support and replace \(f\) by
\(\chi f+(1-\chi)(\rho+1)\), using the exhaustion
\(\rho\) of \cref{p104:M01.07}. This preserves
the signs and the boundary derivative and
agrees with \(\rho+1\) outside a compact set,
so every sublevel set is compact. For the
converse, use a regular coordinate \(f\) at
its zero set and retain the half-space \(f\leq0\).
Away from the zero set use ordinary open
coordinates in the negative locus.
\end{proof}

See \cite[Proposition~5.47, Theorem~5.48, and Problem~5-23]{lee}.
The extra boundary hypothesis in (a) matters. For
\(F:\HH^2\to\R\) given by \(F(x,s)=s+x^2\), \(0\)
is regular for \(F\) but not for its boundary restriction.
Its zero fiber is the single point \((0,0)\) and
has no codimension-one half-slice near that point.
For a different boundary distinction, \(S^1\)
has empty intrinsic boundary while its ambient
frontier in \(\R^2\) is the entire circle.

\section{Smooth coverings}

\begin{proposition}[The smooth structure on a covering]\label{p104:M03.07}
Let \(p:E\to M\) be a topological covering of a
smooth manifold, with countable fibers. Then \(E\)
has a unique smooth structure in which \(p\) is
a smooth covering, meaning it is a diffeomorphism
on each sheet over a sufficiently small connected
evenly covered neighborhood. The same conclusion
holds if \(E\) is connected, since the fibers
are then countable. In either case
\(\partial E=p^{-1}(\partial M)\). If a continuous
map \(\tilde f:X\to E\) lifts a smooth map
\(f:X\to M\), then \(\tilde f\) is smooth.
\end{proposition}

\begin{proof}
Points with distinct images under \(p\) can be
separated by inverse images of disjoint open
neighborhoods, and points in the same fiber
are separated by distinct sheets over a common
evenly covered neighborhood. Thus \(E\) is
Hausdorff. Choose a countable cover by connected
evenly covered coordinate neighborhoods \(U_i\)
in \(M\). If the fibers are countable, there
are countably many sheets above each \(U_i\),
and a countable basis of each of these sheets
gives a countable basis of \(E\).

For the connected case begin with any one sheet
in this cover. A sheet above \(U_i\) meets
at most countably many sheets above any fixed
\(U_j\), since the intersections form disjoint
nonempty open subsets of a second-countable
space. It therefore meets at most countably
many sheets in the whole cover. By induction
only countably many sheets can be reached in
finitely many steps of intersecting sheets.
The union of all such reachable sheets is
open. Its complement is the union of the
remaining sheets, also open, because a
remaining sheet cannot meet a reachable
one. Connectedness implies that the reachable
union is all of \(E\), proving second
countability and countability of the discrete
fibers. If the cover is empty, use the
empty atlas.

Pull back each chart of \(U_i\) along the
homeomorphism on each sheet. Transitions are
base chart transitions restricted to open
overlaps, so the lifted atlas is smooth. It
makes \(p\) a local diffeomorphism and
identifies the boundary from the face of
the base chart. Any structure in which
\(p\) is a local diffeomorphism must include
these charts, proving uniqueness. Finally,
around a point of \(X\) shrink using
continuity so that \(\tilde f\) takes values
in a single sheet. There it is the local
smooth inverse of \(p\) composed with \(f\).
\end{proof}

Consequently the topological path and homotopy lifting
theorems for coverings give smooth lifts whenever
their base maps are smooth, including for boundary
charts. A simply connected topological covering, when
provided by the usual covering existence theorem,
therefore has a canonical smooth structure. For
example the angle map \(\R^k\to(S^1)^k\) and
the antipodal quotient \(S^n\to\R P^n\) from
\cref{p104:M01.03} are smooth coverings, as is seen
on their lifted coordinate neighborhoods.
See \cite[pp.~91--95]{lee}.
\chapter{Sard, Whitney, and transversality}

\section{Critical values}

In this chapter nullity on a manifold means Lebesgue nullity in every chart, using the ambient dimension of the chart even at a boundary. A smooth coordinate change, restricted to small compact Euclidean balls on which it extends, is Lipschitz. A Lipschitz map in the same dimension sends a null set to a null set: replace covering cubes of total volume less than \(\epsilon\) by image cubes whose total volume is at most a constant times \(\epsilon\). Countably many such neighborhoods suffice. This makes the convention independent of the charts. In dimension zero we use counting measure.

\begin{theorem}[Smooth Sard theorem]\label{p104:M04.01}
For a smooth map \(f:M^m\to N^n\), with either manifold allowed to have boundary, its critical values form a null set. A point is critical when \(\dd f_p\) is not surjective. If \(m<n\), the entire image is null. In positive-dimensional targets the regular values are dense. See \cite[Theorem 6.10 and Corollaries 6.11--6.12]{lee} and the corrections to pp.\ 129--131 in \cite{leeerrata}.
\end{theorem}
\begin{proof}
In the Euclidean case let \(f:U\subset\R^m\to\R^n\), with \(U\) open. When \(n=0\) there are no critical points. When \(m=0<n\) the image consists of at most one point. Assume the result in source dimensions below \(m>0\), for every target dimension. Let \(C\) be the critical set and let \(B_j\) be the points where every partial derivative of each component of orders \(1,\ldots,j\) vanishes.

At a point of \(C\setminus B_1\) some first derivative is nonzero. Here \(n\geq2\), since a nonzero derivative into \(\R\) is surjective. Permute coordinates and use the inverse function theorem with the first component of \(f\) as the first source coordinate. In these coordinates \(f(t,x)=(t,g(t,x))\). Its derivative has rank \(1+\operatorname{rank}D_xg(t,x)\). The slice of its critical image at each \(t\) is therefore contained in the critical image of the smooth map \(g(t,\cdot)\) on the open slice, null by induction in dimension \(m-1\). To apply Fubini take a countable exhaustion by compact subsets of the coordinate neighborhood and on each use the image of its intersection with the closed critical set. Each image is compact, with null slices, hence null. Countably many of the rank neighborhoods cover \(C\setminus B_1\).

Near a point of \(B_j\setminus B_{j+1}\), some partial derivative \(h\) of order \(j\) has nonzero differential. The set \(B_j\) there lies in the smooth hypersurface \(h=0\). Choose the neighborhood small enough that this hypersurface has a chart onto an open subset of \(\R^{m-1}\). At each point of \(B_j\) the first derivative of \(f\) is zero, so the corresponding point is critical for the restriction to this chart. The induction hypothesis gives a null image in each such neighborhood. Second countability gives a countable subcover of \(B_j\setminus B_{j+1}\), proving nullity of its image.

Choose \(k\) with \((k+1)n>m\). Cover \(U\) by closed cubes with compact neighborhoods in \(U\). In one such cube subdivide into \(O(N^m)\) subcubes of sidelength \(O(N^{-1})\). In each subcube meeting \(B_k\) choose a point \(p\) of the intersection. On the segment from \(p\) to another point \(q\) of the subcube, Taylor's formula and the bounds on derivatives of order \(k+1\) give
\[
 \lVert f(q)-f(p)\rVert\leq C\lVert q-p\rVert^{k+1}
 \leq C'\,N^{-(k+1)}.
\]
The constants are uniform on the large cube. Thus the image of \(B_k\) in that cube is covered by balls of total \(n\)-volume at most \(C''N^{m-n(k+1)}\), which tends to zero. Combining the three parts proves the induction.

For the manifold assertion choose countably many source and target chart pairs with the image in the target chart. Extend the coordinate map near each source boundary point to a smooth map on a Euclidean open neighborhood. The extension has the same derivative at the point by the density of the model interior. Apply the Euclidean theorem to each extension and restrict its critical image to the target model. A countable refinement suffices. The zero-dimensional source case also follows from second countability of the discrete source. Every point is critical if \(m<n\), and every nonempty relatively open set in the positive-dimensional target chart has positive volume, proving the last two assertions.
\end{proof}

\section{Euclidean embeddings and immersions}

\begin{lemma}[Good projection directions]\label{p104:good-directions}
Let \(e:M^n\to\R^q\) be an injective immersion. If \(q>2n+1\), the projections along a dense set of lines in \(\R^q\) are injective immersions on \(M\). For an immersion alone, it is enough that \(q>2n\). Directions in an arbitrary prescribed proper projective linear subspace can also be excluded.
\end{lemma}
\begin{proof}
An excluded secant is a line spanned by \(e(p)-e(p')\) with \(p\ne p'\). The directions are the image of the smooth map \((p,p')\mapsto[e(p)-e(p')]\) to \(\mathbb{RP}^{q-1}\). When \(M\) has boundary cover the ordered pairs by \((M\times\operatorname{Int}M)\setminus\Delta\) and \((M\times\partial M)\setminus\Delta\), of dimensions \(2n\) and \(2n-1\), respectively. Each is a manifold with at most ordinary boundary. For tangent lines use the map from the unit sphere bundle of \(TM\) to projective space, \((p,v)\mapsto[\dd e_pv]\). It is smooth, its domain has dimension \(2n-1\), and it is empty when \(n=0\). A smooth bundle metric for this purpose is already provided by \(\langle\dd e_pv,\dd e_pw\rangle\); local smooth Gram--Schmidt gives the bundle charts. By \cref{p104:M04.01}, each relevant image is null when its dimension is strictly less than \(q-1\). A proper projective linear subspace is also null. A line outside the exclusions kills neither a nonzero secant nor a nonzero tangent vector when both are required.
\end{proof}

\begin{theorem}[Proper weak Whitney embedding]\label{p104:M04.02}
Every smooth \(n\)-manifold, possibly with boundary, admits a proper smooth embedding in \(\R^{2n+1}\). See \cite[Lemmas 6.13--6.14 and Theorem 6.15]{lee}, with the boundary and exhaustion corrections in \cite{leeerrata}.
\end{theorem}
\begin{proof}
For a compact subset \(K\subset M\), choose finitely many charts \((U_i,x_i)\) and smooth cutoffs \(b_i\) supported in \(U_i\) with the interiors of their plateaus \(b_i=1\) covering \(K\). The map with blocks \(b_i\) and \(b_ix_i\), with the products extended by zero, is an injective immersion on a neighborhood of \(K\). In fact at \(p\) there some \(b_i=1\) near \(p\). Equality of the scalar and vector blocks for \(p,p'\) implies \(p'\in U_i\) with \(x_i(p')=x_i(p)\). The derivative of the same vector block at \(p\) is the coordinate isomorphism on \(T_pM\). This proves the claim even at a boundary point. Use \cref{p104:good-directions} repeatedly, or pad with zeros, to put this injective immersion on the neighborhood in \(\R^{2n+1}\).

Choose a nonnegative smooth proper exhaustion \(a\) as in \cref{p104:M01.07}. For \(j=0,1,\ldots\) let
\[
 K_j=a^{-1}[j,j+1],
 \qquad U_j=a^{-1}(j-\tfrac14,j+\tfrac54).
\]
Each \(\overline{U_j}\) is compact. Construct an injective immersion \(e_j\) into \(\R^{2n+1}\) on a neighborhood of \(\overline{U_j}\) by the preceding paragraph. Let \(\chi_j\) equal one near \(K_j\) and have support in \(U_j\); empty \(K_j\) may be assigned \(\chi_j=0\). The \(U_j\) of either parity are disjoint and the full family is locally finite. The smooth map
\[
 E(p)=\left(\sum_{j\ {\rm even}}\chi_j(p)e_j(p),
             \sum_{j\ {\rm odd}}\chi_j(p)e_j(p),a(p)\right)
\]
is proper because of the last coordinate. If \(E(p)=E(p')\), choose \(j\) with \(p\in K_j\); equality of the last coordinates gives \(p'\in K_j\). Their parity blocks then read \(e_j(p),e_j(p')\), which imply \(p=p'\). The same block is an immersion near \(p\). Properness and \cref{p104:M03.02} show that \(E\) is an embedding. This construction uses only open neighborhoods of the compact exhaustion bands, which need not themselves be smooth domains.

To lower the dimension while retaining properness, first replace a finite-dimensional embedding \(E_0\) by \((G\circ E_0,a)\), where \(G:\R^q\to B(0,1)\) is the diffeomorphism \(z\mapsto z/\sqrt{1+\lVert z\rVert^2}\). Its image lies a bounded distance from a line \(S\), the last coordinate axis, and the new map is proper. Suppose such a proper embedding lies in \(\R^r\) with \(r>2n+1\). Project along a good line \(L\ne S\) using \cref{p104:good-directions}; take the orthogonal projection \(P\) to \(L^\perp\). The projected image lies within a bounded distance of the line \(P(S)\). If \(P(x)\) lies in a compact set, then \(\operatorname{dist}(x,L)=\lVert P(x)\rVert\) is bounded. On the unit sphere the continuous function
\[
 u\longmapsto \max(\operatorname{dist}(u,L),\operatorname{dist}(u,S))
\]
has a strictly positive minimum, since the two lines intersect only in zero. Therefore the intersection of the two bounded-radius tubes is bounded. The inverse image of a compact set under the projected map is closed in the inverse image of a sufficiently large closed ball under the original proper map. This proves properness at each projection step. The projected injective immersion is again an embedding. Iterate, padding with zeros only if needed.
\end{proof}

\begin{theorem}[Weak immersion and compact approximation by embeddings]\label{p104:M04.03}
Every smooth \(n\)-manifold, with or without boundary, immerses in \(\R^{2n}\). For compact \(M^n\) and \(q\geq 2n+1\), every smooth \(f:M\to\R^q\) is uniformly approximable by smooth embeddings into \(\R^q\). See \cite[Corollary 6.17 and Theorem 6.18]{lee}.
\end{theorem}
\begin{proof}
Apply the tangent part of \cref{p104:good-directions} to the weak embedding into \(\R^{2n+1}\). For the approximation, start from the embedding \((f,e)\) into \(\R^q\times\R^{2n+1}\). Project off the added coordinates one at a time. At each step use a good line arbitrarily near the last coordinate axis and the projection onto the coordinate hyperplane along that line. The projected map stays an injective immersion until the ambient dimension reaches \(q\). Composition of the finitely many projections can be made arbitrarily close in operator norm to the original coordinate projection: each projection varies continuously with its line near the corresponding coordinate axis, and the good choices are dense at every step. Since \((f,e)(M)\) is bounded, this gives the desired uniform approximation. Compactness makes each injective immersion an embedding.
\end{proof}

The stronger Whitney bounds for boundaryless manifolds are embedding into \(\R^{2n}\) when \(n>0\) and immersion into \(\R^{2n-1}\) when \(n>1\). These are cited bounds here, with the stronger proofs assigned to differential topology \cite[Theorems 6.19--6.20]{lee}.

\section{Tubular neighborhoods and smooth approximation}

\begin{lemma}[Shrinking around an embedded image]\label{p104:shrink-along-section}
Let \(F:U\to N\) be continuous, and let \(S\subset U\) be a subspace on which \(F\) is a homeomorphism onto a locally closed subset of the manifold \(N\). Suppose \(F\) is injective on some open neighborhood of each point of \(S\) in \(U\). Then \(F\) is injective on an open neighborhood of \(S\).
\end{lemma}
\begin{proof}
Restrict the target to an open set \(O\) in which \(F(S)\) is closed and restrict the source to \(F^{-1}(O)\). For every \(s\in S\), choose an injectivity neighborhood \(A_s\) and an open \(O_s\subset O\) around \(F(s)\) such that \(F^{-1}(O_s)\cap S\subset A_s\). This is possible because \(F|_S\) is a homeomorphism. Add \(O\setminus F(S)\) to the target cover and assign to it the empty injectivity set. By regularity, around each point of \(O\) choose a neighborhood whose closure in \(O\) lies in a member of this cover. By paracompactness, the cover by these smaller neighborhoods has a locally finite open refinement \(W_i\). Each \(\overline W_i\) lies in the closure of one smaller neighborhood and thus in its assigned original member. If \(A_i\) is the corresponding injectivity set, remove from \(F^{-1}(O)\) the union of
\[
 F^{-1}(\overline W_i)\setminus A_i .
\]
The closures remain locally finite, since an open neighborhood meeting a closure meets the corresponding open set. Their preimages are locally finite by continuity, so the displayed union is closed. It misses \(S\). For two surviving points with common image \(y\), choose \(i\) with \(y\in W_i\). Both points belong to \(A_i\), so they coincide.
\end{proof}

\begin{theorem}[Euclidean tubular neighborhood]\label{p104:M04.04}
If \(S\) is a boundaryless embedded submanifold of \(\R^q\), the normal bundle has an open neighborhood of the zero section on which \((p,v)\mapsto p+v\) is a diffeomorphism onto an open neighborhood of \(S\). The neighborhood can be chosen fiberwise star-shaped, giving a smooth retraction and a deformation to \(S\). See \cite[Theorem 6.24 and Proposition 6.25]{lee}.
\end{theorem}
\begin{proof}
The normal spaces are the orthogonal complements of the tangent spaces. A smooth tangent frame on a chart can be extended near a point to a full ambient frame and orthogonalized by Gram--Schmidt. This gives local smooth bundle charts for the normal spaces. At \((p,0)\) the tangent space of the total space canonically splits into the tangent of the zero section and the vertical normal space. The derivative of addition is \((w,u)\mapsto w+u\), an isomorphism, so addition is a local diffeomorphism near the zero section. The embedded slice charts imply that \(S\) is locally closed. Apply \cref{p104:shrink-along-section}, restricting also to the open local-diffeomorphism locus. If \(V\) is the resulting open neighborhood, cover \(S\) by chart neighborhoods on each of which all normal vectors with norm less than a constant \(a_i>0\) lie in \(V\). A subordinate locally finite partition \(\rho_i\) gives a positive smooth radius \(\epsilon(p)=\sum_i\rho_i(p)a_i\). At each \(p\), this is at most the largest of the active admissible radii. The open disk bundle of radius \(\epsilon\) thus lies in \(V\). Its image is open by local invertibility, and its inverse is smooth locally and hence globally. Projection to the base gives the retraction, and radial contraction of normal vectors gives the deformation.
\end{proof}

\begin{theorem}[Relative Whitney approximation]\label{p104:M04.05}
Let \(F:M\to\R^q\) be continuous and let \(A\subset M\) be closed. Assume that at each point of \(A\), \(F|_A\) has a smooth local extension. For every positive continuous \(\epsilon:M\to\R\) there is a smooth \(G\) with \(G|_A=F|_A\) and \(\lVert G(p)-F(p)\rVert<\epsilon(p)\). There is a homotopy from \(F\) to \(G\) fixed on \(A\). For a boundaryless manifold target the same relative homotopy conclusion holds, and the homotopy can be kept inside any prescribed open neighborhood of the graph of \(F\). The source may have boundary. Compare \cite[Theorems 6.21, 6.26--6.27]{lee} and the correction to p.\ 45 in \cite{leeerrata}.
\end{theorem}
\begin{proof}
Near a point of \(A\) use the given extension, shrinking its domain until its distance from \(F\) is less than \(\epsilon\); continuity and equality at that point make this possible. Away from \(A\) use constants \(F(p)\) on small neighborhoods disjoint from \(A\) satisfying the same error bound. Average these smooth functions with a subordinate smooth partition, extending the weighted terms by zero. All nonzero terms at points of \(A\) have the prescribed value. The pointwise convex combination gives the error estimate. The straight homotopy is fixed on \(A\).

For a boundaryless target embed it properly into Euclidean space by \cref{p104:M04.02}, and let \(r\) be the retraction of a tubular open neighborhood. Given an open neighborhood \(\mathcal O\subset M\times N\) of the graph, for each \(p\) choose a source neighborhood \(V_i\) and a number \(\delta_i>0\) so that for \(x\in V_i\) the ball of radius \(\delta_i\) about the embedded \(F(x)\) is inside the tubular domain, and its image under \(r\) lies in \(\mathcal O_x\). This follows by continuity of \(F\) and of \((x,z)\mapsto (x,r(z))\), using product neighborhoods at the graph points. Patch the radii with a partition as in the preceding proof to obtain a positive continuous tolerance. Approximate the embedded \(F\) within that tolerance. Compose the entire straight homotopy with \(r\). The result stays in \(\mathcal O\), is fixed on \(A\), and ends at a smooth map. If \(F\) is smooth on a neighborhood of \(A\), the resulting homotopy is smooth on a neighborhood of that subset times the interval.
\end{proof}

\section{Transverse inverse images and parameters}

\begin{theorem}[Transverse inverse image]\label{p104:M04.06}
Let \(S\subset N\) be a boundaryless embedded submanifold of a boundaryless manifold, and let \(f:M\to N\) be smooth. Suppose
\[
 \dd f_p(T_pM)+T_{f(p)}S=T_{f(p)}N
 \qquad(f(p)\in S).
\]
If \(M\) has boundary, assume also the same condition with \(T_p\partial M\) at points of \(\partial M\) mapping into \(S\). Then \(f^{-1}(S)\) is an embedded submanifold of codimension \(\operatorname{codim}_N S\) wherever nonempty, its boundary is \(f^{-1}(S)\cap\partial M\), and its tangent space is \((\dd f_p)^{-1}(T_{f(p)}S)\). See \cite[Theorem 6.30 and Problem 5-23]{lee}.
\end{theorem}
\begin{proof}
Near \(f(p)\), write \(S\) as the zero set of the normal coordinate projection \(g\) of a slice chart. Then \(\ker\dd g=T_{f(p)}S\), so the displayed condition is precisely surjectivity of \(\dd(g\circ f)_p\). At a source boundary point the additional assumption says that its restriction to \(T_p\partial M\) is surjective. Apply the interior or boundary regular-level result from \cref{p104:M03.05,p104:M03.06}. The local tangent kernel is the claimed inverse image. Since the structures come from ambient slice charts, these descriptions agree on overlaps.
\end{proof}

\Needspace{17\baselineskip}
\subsection{Families and perturbations}

\begin{theorem}[Parametric transversality]\label{p104:M04.07}
Let \(P\) be a boundaryless parameter manifold and let \(F:P\times M\to N\) be smooth and transverse to \(S\) as in \cref{p104:M04.06}. If \(M\) has boundary assume transversality also on \(P\times\partial M\). For all parameters outside a null set, \(F_a=F(a,\cdot)\) is transverse to \(S\), including on \(\partial M\) when present. This is the projection proof of \cite[Theorem 6.35]{lee}, with the linear inclusion at p.\ 146 as corrected in \cite{leeerrata}.
\end{theorem}
\begin{proof}
The inverse image \(W=F^{-1}(S)\) is a submanifold with the stated boundary. At \((a,p)\in W\) put \(Q=T_{F(a,p)}N/T_{F(a,p)}S\), and denote the two derivative blocks modulo \(T S\) by \(A:T_pM\to Q\) and \(B:T_aP\to Q\). Transversality of \(F\) is surjectivity of \((A,B)\). The tangent of \(W\) consists of \((b,v)\) with \(B b+A v=0\). Its projection to \(T_aP\) is onto if and only if \(\operatorname{im}B\subseteq\operatorname{im}A\); together with surjectivity of \((A,B)\) this is equivalent to surjectivity of \(A\). The same argument on \(\partial W\) uses \(T_p\partial M\). Apply Sard to \(W\to P\) and, when needed, to \(\partial W\to P\), and exclude the union of the two null sets.
\end{proof}

For example, for Euclidean targets, \(F(a,p)=f(p)+a\) has surjective parameter derivative and so does its boundary restriction. Arbitrarily small choices of \(a\) give transverse perturbations homotopic to \(f\). If \(M\) is boundaryless and a smooth homotopy \(H:[0,1]\times M\to\R^n\) has transverse endpoints, apply the theorem to \(H(t,p)+t(1-t)a\). The full parameter derivative is surjective for \(0<t<1\), while the endpoint restrictions are already transverse. This gives a transverse homotopy with unchanged endpoints, whose inverse image has boundary exactly the two endpoint inverse images. This argument uses a finite-dimensional family for which the parameter derivative is controlled.
\chapter{Vector bundles, covectors, and tensors}

\section{Bundles assembled from local data}

A smooth real vector bundle of rank \(r\) over \(M\) is a smooth map \(\pi:E\to M\) with real vector-space fibers and smooth local trivializations \(E|_U\cong U\times\R^r\), linear on fibers. On an overlap write the change from the \(i\)-coordinates to the \(j\)-coordinates as \((p,v)\mapsto(p,g_{ji}(p)v)\). The matrix \(g_{ji}\) is a smooth map to \(\mathrm{GL}_r(\R)\), and \(g_{ki}=g_{kj}g_{ji}\) on triple overlaps. A rank-zero bundle has only zero fibers.

\begin{theorem}[Gluing from transition functions]\label{p104:M05.01}
Given a cover of \(M\) and smooth transition functions with \(g_{ii}=I\) and the cocycle identities, there is a vector bundle with those transitions. Two realizations are uniquely isomorphic by the map preserving the given trivializations. This is the construction underlying \cite[Chapter 10 and Problem 10-6]{lee} and \cite[\S12]{tu}.
\end{theorem}
\begin{proof}
On \(\coprod_i U_i\times\R^r\) set \((p,v,i)\sim(p,g_{ji}(p)v,j)\) whenever \(p\in U_i\cap U_j\). The identities give reflexivity, symmetry and transitivity. In the quotient topology each \(U_i\times\R^r\) maps homeomorphically onto the full inverse image of \(U_i\). To verify openness for an open subset of this chart, take its inverse image in each other chart by the given continuous transition homeomorphism. The result is open in the disjoint union and hence gives an open set of the quotient. Points over different bases separate by pullback of disjoint base neighborhoods; points over the same base separate in one common trivialization. A countable subcover of base trivialization domains yields a countable collection of countable bases upstairs. Smooth transitions give the manifold structure on the total space, including relative half-space charts over \(\partial M\). Fiber operations and the zero section are smooth in these charts. The local coordinate-preserving maps between any two realizations glue because they give the same transition rule.
\end{proof}

For the tangent bundle, the transitions of \cref{p104:M02.04} are the derivatives of the base coordinate changes. The cocycle is exactly the chain rule; smoothness of these matrices is the reason the construction gives a smooth bundle. A frame is an ordered list of sections giving a basis on every fiber in its domain.

\begin{proposition}[Sections, frames, and extension]\label{p104:M05.02}
In a smooth local frame a section is smooth exactly when its coefficient functions are smooth. A rank-\(r\) bundle is trivial exactly when it has a global smooth frame. A section prescribed on a closed subset \(A\subset M\) and smoothly extendable near each point of \(A\) admits a global smooth extension. The extension can have support in any given open \(U\supset A\) if local extensions are taken inside \(U\). Compare \cite[pp.\ 255--261]{lee} and \cite[\S\S12.4--12.5]{tu}.
\end{proposition}
\begin{proof}
In a trivialization the columns of a frame give an invertible matrix \(B(p)\). Smoothness of \(B\) and its inverse proves the local coefficient test and the characterization of a trivialization by a frame. To extend the section, cover \(M\) by the local extension domains in \(U\) and \(M\setminus A\), assigning the zero section to the latter. Multiply the local sections by a subordinate smooth partition and extend each weighted section by zero. The sum is locally finite and all contributions at a point of \(A\) have the same prescribed value before weighting. The support of the result is in the union of the supports of the weights for the domains inside \(U\). This union is closed by local finiteness and is contained in \(U\). In particular multiplying a local frame member by a cutoff gives a global section agreeing with it near a chosen point.

For a section of a restricted bundle over a closed embedded submanifold, local extensions also exist at boundary points. Choose an ambient trivialization and local Euclidean or half-space coordinates on both manifolds. Extend the coordinate inclusion and the given section coefficients to Euclidean neighborhoods of the source point. The derivative of the inclusion has full source rank, which persists on a smaller Euclidean neighborhood. The Euclidean immersion normal form provides ambient coordinates straightening the extended local image. Extend the coefficients in those coordinates by making them constant in the normal variables and restrict the result to the actual ambient model. Because the original inclusion is an embedding, the ambient neighborhood can be chosen so that all the submanifold points in it belong to the chosen source chart. The resulting section therefore agrees with the original data on the submanifold near the point.
\end{proof}

\begin{theorem}[Kernels, images, and quotients]\label{p104:M05.03}
A smooth bundle morphism \(A:E\to F\) over \(\id_M\), of constant rank, has smooth kernel and image subbundles. If \(K\subset E\) is a smooth subbundle, the fiberwise quotients \(E_p/K_p\) give a smooth quotient bundle. These assertions also apply componentwise when the ranks vary between components. See \cite[pp.\ 261--268]{lee}.
\end{theorem}
\begin{proof}
At a point of rank \(k\), permute local frame coordinates so that the \(k\times k\) upper-left minor \(B\) is invertible. Write the matrix as
\(
 \begin{psmallmatrix}B&C\\D&H\end{psmallmatrix}.
\)
Smooth block row and column operations reduce it to diagonal blocks \(B\) and \(H-DB^{-1}C\). The second is zero by the rank hypothesis. Multiplying the first block by \(B^{-1}\) gives a matrix with blocks \(I_k,0\). Thus the local smooth changes of frame identify kernel and image with constant coordinate subspaces. For \(k=0\) the map is zero in the chosen neighborhood. For the quotient, extend a local frame of \(K\) to one of \(E\), using the persistence of linear independence near each point. The remaining vectors project to a basis of \(E_p/K_p\). Transition matrices of adapted frames have the block preserving \(K\), and their induced matrices on the quotient are smooth and invertible. The preceding gluing theorem gives the quotient. Its projection is smooth, fiberwise surjective, and has the usual factorization property for morphisms vanishing on \(K\); smoothness of the factorization follows in an adapted frame.
\end{proof}

\section{Functorial constructions and metrics}

\begin{proposition}[Pullbacks and algebraic operations]\label{p104:M05.04}
If \(f:P\to M\) is smooth, the fibers \(E_{f(p)}\) form \(f^*E\) over \(P\). Pullback along a composite is canonically isomorphic to the iterated pullback. Duals, direct sums, tensor products and exterior powers are smooth bundles. The normal bundle of an embedded submanifold \(S\subset M\) is \(TM|_S/TS\). The dual of \(TM\) is the cotangent bundle \(T^*M\). See \cite[Chapters 10--12]{lee} and \cite[\S\S17--18]{tu}.
\end{proposition}
\begin{proof}
For \(f^*E\) pull each chart and matrix \(g_{ji}\) back along \(f\). The cocycle is unchanged, and the fiber product description \(\{(p,e):f(p)=\pi(e)\}\) gives the canonical smooth map to \(E\). For a composite the two models agree fiberwise with the identity on the indicated copy of \(E_{f(g(p))}\), and the local matrices are equal. Dual coordinates change by \(g_{ji}^{-\mathsf T}\), direct sums by block diagonal matrices, tensors by products of matrix entries, and exterior powers by minors. Each is compatible with matrix composition, so the cocycles are preserved. Smoothness follows from the indicated polynomial operations or smooth matrix inversion. For the inclusion \(i:S\to M\), the derivative gives a smooth fiberwise injection \(TS\to i^*TM\) over \(\id_S\), also at boundary points. Its image is a smooth subbundle by \cref{p104:M05.03}, which also constructs the normal quotient. The same local conventions give \(\Lambda^0 E=M\times\R\) and zero bundles in exterior degrees greater than the fiber rank.
\end{proof}

\begin{theorem}[Metrics and splitting]\label{p104:M05.05}
Every smooth real vector bundle admits a smooth positive-definite fiber metric. If \(K\subset E\) is a smooth subbundle, its metric orthogonal complement is smooth and gives \(E=K\oplus K^\perp\). Consequently short exact sequences of smooth vector bundles split. A fiber metric identifies a bundle with its dual by \(v\mapsto \langle v,\cdot\rangle\). See \cite[Chapter 13]{lee}.
\end{theorem}
\begin{proof}
In each trivialization take the standard inner product and average with a subordinate partition of unity. In each fiber the result is a finite convex combination of positive-definite forms and is therefore positive definite. For a local frame \(k_1,\ldots,k_s\) of \(K\), let \(G_{ij}=\langle k_i,k_j\rangle\). The orthogonal projection is
\[
 v\longmapsto \sum_{i,j}k_i(G^{-1})_{ij}\langle k_j,v\rangle .
\]
It is smooth by the matrix inverse formula. Its kernel is \(K^\perp\) and is a subbundle by \cref{p104:M05.03}; the local formula also gives the smooth projections for the direct sum. In a short exact sequence identify the injection with its image and restrict the quotient map to the orthogonal complement. The restriction is a smooth fiberwise isomorphism, whose inverse is smooth in local matrices. The metric map to the dual and its inverse are likewise represented by smooth invertible Gram matrices. In particular on \(TM\) it gives the usual identification with \(T^*M\). Under a change of tangent basis \(B\), the Gram matrix becomes \(B^{\mathsf T}GB\), so its positive square-root determinant changes by the factor \(\lvert\det B\rvert\). This is the transformation used for volume densities.
\end{proof}

\begin{proposition}[Differentials and covector pullback]\label{p104:M05.06}
A one-form on \(M\) is a smooth section of \(T^*M\). In a coordinate chart \(x=(x^1,\ldots,x^n)\), the covectors \(\dd x^i\) are the dual frame to the coordinate vector fields. The differential of a smooth function \(u\) is a smooth one-form with local coefficients \(\partial_i u\). A smooth map \(f:P\to M\) pulls one-forms back by
\[
 (f^*\alpha)_p(v)=\alpha_{f(p)}(\dd f_p v),
 \qquad \dd(u\circ f)=f^*(\dd u).
\]
The construction respects compositions and restriction to submanifolds. See \cite[pp.\ 272--287]{lee} and \cite[\S17]{tu}.
\end{proposition}
\begin{proof}
The coordinate frame is defined by the tangent chart. Its dual is the differential of the coordinate map componentwise, by the coordinate formula for the tangent map. The displayed formula is application of the dual of \(\dd f_p\) to the pulled-back section; in local coordinates its coefficients are products of smooth functions with entries of \(Df\). The chain rule gives the equation for \(\dd u\) and functoriality. Restriction as a differential form uses the differential of the inclusion, so it evaluates only on tangent vectors to the submanifold, unlike restriction merely as an ambient covector over points of that submanifold.
\end{proof}

\section{Tensors and the abstract tubular theorem}

\begin{proposition}[Tensor fields and their transformations]\label{p104:M05.07}
Let \(T^r_sM=(TM)^{\otimes r}\otimes(T^*M)^{\otimes s}\). Symmetrization and alternation on like slots, tensor products and contraction of paired tangent and cotangent slots define smooth bundle operations. Pullback of a covariant tensor is defined for any smooth map. Pullback of arbitrary mixed tensors is defined along a local diffeomorphism and commutes with tensor products and all permitted contractions. See \cite[Chapter 12]{lee} and \cite[\S3]{tu}.
\end{proposition}
\begin{proof}
A local frame and its dual identify tensor coordinates with arrays indexed by tangent and covector bases. Multilinearity gives the canonical pairing with \(r\) covector arguments and \(s\) vector arguments. The identification is invertible since values on all basis tuples recover the array. Permutations of like slots act by linear permutations of the entries. Averaging them with coefficients \(1/k!\) or \(\operatorname{sgn}(\sigma)/k!\) gives the two projections. The permutation identities and \(\sum_\sigma 1=k!\) show idempotence and the expected symmetry of the image. Contraction is induced by the evaluation map \(\lambda(v)\), so it is basis-independent. In a paired index it sums equal upper and lower indices. All these local formulas are smooth. For a covariant tensor, evaluate it on the images of the vector arguments under \(\dd f_p\). For a mixed tensor, use \(\dd f_p\) to transport the vector arguments and the dual of the inverse linear map \((\dd f_p)^{-1}\) to transport the covector arguments. Equivalently apply the inverse differential to contravariant tensor factors and the dual of the differential to covariant factors. In a contraction the two induced maps cancel by the natural pairing, and tensor products are preserved by multilinearity. Composition follows from the chain rule. A map with noninvertible differential still has a pullback bundle \(f^*(T^r_sM)\), with fibers given by the tensor spaces at \(f(p)\). Transport into \(T^r_sP\) requires the factor maps described above.
\end{proof}

\begin{theorem}[Tubular neighborhoods in abstract ambient manifolds]\label{p104:M05.08}
Let \(S\) be a boundaryless embedded submanifold of a boundaryless \(M\). An open neighborhood of the zero section in \(TM|_S/TS\) is diffeomorphic to an open neighborhood of \(S\) in \(M\), by a map restricting to the inclusion on the zero section. The domain can be chosen fiberwise star-shaped. This follows from the Euclidean construction in \cite[Theorems 6.15 and 6.24 and Proposition 6.25]{lee}.
\end{theorem}
\begin{proof}
By \cref{p104:M04.02} identify \(M\) with an embedded submanifold of \(\R^q\), and let \(r:W\to M\) be its Euclidean tubular retraction from \cref{p104:M04.04}. Identify the quotient normal bundle of \(S\) with the orthogonal complement \(K=(TS)^\perp\) inside \(TM|_S\), using the metric induced from Euclidean space and \cref{p104:M05.05}. Where \(p+v\in W\), set \(F(p,v)=r(p+v)\). On the zero section \(F\) is the inclusion. At \((p,0)\), its differential sends \((w,u)\) to \(\dd r_p(w+u)=w+u\), because \(\dd r_p\) is the identity on \(T_pM\). It is therefore invertible at each zero vector. Restrict to the local-diffeomorphism locus and use \cref{p104:shrink-along-section} for the locally closed embedded image of the zero section. The resulting injective local diffeomorphism has an open image and a smooth inverse. Shrink further to an open disk bundle of smoothly varying radius by the radius argument in \cref{p104:M04.04}. This proves the assertion without a geodesic construction.
\end{proof}
\chapter{Vector fields, flows, and collars}

\section{Fields and brackets}

\begin{proposition}[Vector fields as derivations]\label{p104:M06.01}
Smooth sections \(X\) of \(TM\) correspond to real-linear derivations of \(C^\infty(M)\), via \(Xf(p)=\dd f_pX(p)\). In a chart the coefficients of \(X\) are \(Xx^i\). For a smooth \(f:M\to N\), vector fields \(X\) on \(M\) and \(Y\) on \(N\) are \(f\)-related precisely when \(X(u\circ f)=(Yu)\circ f\) for every smooth \(u\) on \(N\). See \cite[pp.\ 174--185]{lee} and \cite[\S14]{tu}.
\end{proposition}
\begin{proof}
A smooth section gives a derivation by the product rule. Conversely a derivation \(D\) is local: if \(g\) vanishes near \(p\), choose a bump \(b=1\) near \(p\) supported where \(g=0\). Then \(0=D(bg)(p)=Dg(p)\). The same reasoning permits \(D\) to act on locally defined functions by using extensions obtained from cutoffs, with no change near \(p\). Thus \(g\mapsto Dg(p)\) is a point derivation. By \cref{p104:M02.01} it gives a unique \(X(p)\). On any smaller chart with the coordinate functions extended using a cutoff, the coefficients are the smooth functions obtained by applying \(D\) to those extensions. Hence \(X\) is smooth. Finally \(\dd f_pX(p)=Y(f(p))\) implies the related-functions formula by the chain rule. In the reverse direction coordinate functions times cutoffs have differentials spanning \(T^*_{f(p)}N\), so the formula implies equality of the vectors.
\end{proof}

When \(f\) is a diffeomorphism every field pushes forward by \((f_*X)(q)=\dd f_{f^{-1}(q)}X(f^{-1}(q))\). The same statement is local for a local diffeomorphism. A general smooth map may send different points of one fiber to different candidate vectors, so related fields have to be specified.

\begin{proposition}[Lie bracket]\label{p104:M06.02}
The commutator \([X,Y]f=X(Yf)-Y(Xf)\) is a smooth vector field. In coordinates its \(j\)-th coefficient is
\[
 \sum_i\bigl(X^i\partial_iY^j-Y^i\partial_iX^j\bigr).
\]
The bracket is real-bilinear and antisymmetric, satisfies the Jacobi identity and \([X,hY]=h[X,Y]+(Xh)Y\), and preserves related fields under smooth maps. For \(C^r\) fields with integer \(r\geq1\) the coordinate formula yields a \(C^{r-1}\) field. See \cite[pp.\ 185--199]{lee} and \cite[\S14.4]{tu}.
\end{proposition}
\begin{proof}
In the commutator applied to a product the two mixed first-derivative terms occur in both orders and cancel. The result is a derivation, hence a vector field by \cref{p104:M06.01}. Applying it to each coordinate gives the displayed coefficients. Expanding the three double commutators in the Jacobi identity cancels the six possible ordered compositions in pairs. Expanding the product by \(h\) gives the Leibniz identity. If \(X,Y\) are related to \(X',Y'\) under \(f\), apply the two related-functions identities successively to \(u\) and to \(Y'u\) or \(X'u\). Subtraction shows \([X,Y]\) related to \([X',Y']\). The same coordinate expression and chain rule give the finite regularity statement.
\end{proof}

For instance, on a Lie group \(G\), a vector \(v\in T_eG\) gives the smooth left-invariant field \(X_v(g)=\dd(L_g)_e v\), using smoothness of the multiplication. Every left-invariant field is obtained this way. Since all left translations preserve these fields, they preserve their brackets. The bracket on \(T_eG\) is consequently the value at \(e\) of \([X_v,X_w]\). On the open matrix group \(\mathrm{GL}_n(\R)\) the field is \(X_v(g)=gv\), and the coordinate formula gives \([X_v,X_w](g)=g(vw-wv)\).

\section{Maximal flows}

\begin{theorem}[Maximal integral curves and their flow]\label{p104:M06.03}
Let \(M\) be boundaryless and \(X\) a smooth vector field. For each \(p\) there is a unique maximal integral curve \(\gamma_p:I_p\to M\), with \(0\in I_p\), \(\gamma_p(0)=p\), and open interval \(I_p\). The set
\[
 \mathcal D=\{(t,p)\in\R\times M:t\in I_p\}
\]
is open and \(\varphi(t,p)=\gamma_p(t)\) is smooth on \(\mathcal D\). If \(s\in I_p\), then \(I_{\varphi_s(p)}=I_p-s\) and \(\varphi_t(\varphi_s(p))=\varphi_{s+t}(p)\) for \(t+s\in I_p\). In particular each \(\varphi_s\) is a diffeomorphism from its open domain onto that of \(\varphi_{-s}\). Compare \cite[Theorems 9.12--9.13 and Appendix D]{lee}, \cite[\S\S14.2--14.3]{tu}, and the Euclidean flow results of \cite[Chapters 6 and 11]{d13}.
\end{theorem}
\begin{proof}
A chart identifies an integral curve with the ODE for the coefficient map of \(X\). In an open state domain, the Euclidean local theorem gives around each initial point a common time interval and a smooth solution map. The chain rule identifies curves obtained in different charts. Two curves through a given point at a given time agree throughout the intersection of their time intervals: the coincidence set there is closed by Hausdorffness and open by local uniqueness, and contains that time. Gluing all intervals containing zero gives one curve on their union, an open interval, and its construction makes it maximal.

Fix \((t_0,p_0)\in\mathcal D\), first with \(t_0\geq0\). The compact trajectory segment \(\gamma_{p_0}([0,t_0])\) is covered by finitely many open initial-state neighborhoods on which chart solutions are available for uniform positive and negative times. Choose a finite partition \(0=s_0<\cdots<s_k=t_0\) with lengths smaller than the minimum of these time bounds; if \(t_0=0\), use one local solution. At each step choose a neighborhood containing \(\gamma_{p_0}(s_i)\), and compose its local time map with the preceding steps. By continuity shrink the starting neighborhood at \(p_0\) finitely many times so every composition stays in the next initial-state neighborhood. One more local map at the end uses time \(t-t_0\) for \(t\) near \(t_0\). The concatenated curve agrees with the unique maximal solution by local uniqueness. This proves openness and smoothness near \((t_0,p_0)\). For \(t_0<0\) traverse the segment in decreasing time. The Euclidean maps for each differentiability order can be taken on this same canonical solution, so smoothness follows at every finite order.

Translation of a maximal solution by \(s\) gives an integral curve through \(\varphi_s(p)\) on \(I_p-s\). If it could be extended, translating back would extend the original one. This proves equality of lifespans and the flow law. Taking \(t=-s\) in both directions yields the two inverse smooth time maps.
\end{proof}

\begin{corollary}[Escape and completeness]\label{p104:M06.04}
At a finite endpoint of a maximal interval, the integral curve eventually leaves each compact subset of \(M\). A compactly supported vector field on a boundaryless manifold is complete. In particular every smooth vector field on a compact boundaryless manifold is complete. See \cite[pp.\ 213--217]{lee}.
\end{corollary}
\begin{proof}
If times approach a finite endpoint while the states still visit a compact \(K\), finitely many uniform-time solution neighborhoods of \(K\) give a single positive time bound. Restarting from a visit within that bound of the endpoint extends the trajectory past the endpoint, a contradiction. The argument is symmetric for the left endpoint. For a compactly supported field a curve that meets the complement of the support is constant by uniqueness through that point. Otherwise its image is contained in the compact support. In either case the escape assertion excludes a finite endpoint. On a compact manifold the support is compact.
\end{proof}

\section{Flow boxes and the boundary}

\begin{theorem}[Flowout and collar]\label{p104:M06.05}
At a point where a smooth field is nonzero on a boundaryless manifold, it is the first coordinate field in local flow coordinates. More generally a boundaryless smooth embedded hypersurface transverse to the field has a flowout on a neighborhood of its zero-time section. If \(M\) has boundary, it admits a smooth collar \(c:\partial M\times[0,1)\to M\) onto an open neighborhood of \(\partial M\), with \(c(q,0)=q\). The latter uses the one-sided flow of an inward-pointing field. Compare \cite[Theorems 9.20, 9.24--9.25 and Problem 9-11]{lee}.
\end{theorem}
\begin{proof}
For a transverse hypersurface \(S\), use \(F(t,q)=\varphi_t(q)\). At \((0,q)\) the differential is \((a,v)\mapsto aX(q)+v\), an isomorphism by transversality. The inverse function theorem gives local flow coordinates, and the flow law gives \(F_*\partial_t=X\). Around a single nonzero vector choose a small coordinate hypersurface transverse to it, which proves the flow-box statement. On the whole zero-time section, \(F|_{\{0\}\times S}\) is the embedded inclusion. By \cref{p104:shrink-along-section}, restrict its local-diffeomorphism locus to an open neighborhood on which it is injective. This gives a flowout by a diffeomorphism onto its image, without a common time width on noncompact \(S\).

For the boundary choose a smooth field \(X\) pointing strictly inward at every boundary point. In a boundary chart use \(\partial_{x_n}\), and take a subordinate partition of unity for the boundary chart domains and the interior, assigning the zero field to the interior. At a boundary point, the inward condition is positivity in the inward half of the one-dimensional quotient \(T_pM/T_p\partial M\). The positive local contributions preserve this condition. Extend the coefficient functions of \(X\) near a boundary point to a Euclidean open set and choose a smaller neighborhood where the extended normal component is strictly positive. By continuity of the Euclidean local flow, shrink both the initial-state neighborhood and its uniform time interval so that every trajectory used stays inside that smaller neighborhood. For boundary initial states the normal coordinate is then strictly increasing from zero. The local curve lies in \(M\) for nonnegative time and leaves it for negative time. On an overlap, the two local trajectories from a given boundary point agree throughout their common nonnegative time interval. The coincidence set is closed by Hausdorffness and relatively open by local Euclidean uniqueness, using a common chart and one extension near each coincident point. It contains time zero. Taking the union of these neighborhoods thus gives a smooth map \(F(t,q)\) on a relative open neighborhood of \(\{0\}\times\partial M\) in \([0,\infty)\times\partial M\).

At \((0,q)\) the extended flow map has invertible differential for the same transversality reason. Its normal coordinate has the sign of \(t\) near the zero-time section, so restriction of the Euclidean local inverse gives a relative local diffeomorphism of half-space models. Since the boundary is closed, \cref{p104:shrink-along-section} gives an injective relative local diffeomorphism on an open neighborhood of the zero-time section. Choose a positive smooth function \(\epsilon(q)\) such that \(0\leq t<\epsilon(q)\) lies in that neighborhood: local constant lower bounds for the width patch by a partition of unity, as in \cref{p104:M04.04}. Now set \(c(q,t)=F(t\epsilon(q),q)\) for \(0\leq t<1\). This is a diffeomorphism onto an open neighborhood of the boundary. The construction works componentwise and the empty-boundary case gives the empty collar.
\end{proof}

\begin{proposition}[Boundary-tangent flows and smoothing homotopies]\label{p104:M06.06}
If \(X\) is tangent to \(\partial M\), its maximal flow exists for two-sided local times with open domain relative to \(\R\times M\). It preserves the boundary and the interior, and satisfies the laws of \cref{p104:M06.03}. The inclusion \(\operatorname{Int}M\hookrightarrow M\) is a smooth homotopy equivalence. Every continuous map into a manifold with boundary is homotopic to a smooth map, and a continuous homotopy between smooth maps can be replaced by a smooth homotopy with the same endpoints. Compare \cite[Theorems 9.26--9.29 and 9.34--9.35]{lee}.
\end{proposition}
\begin{proof}
Extend the coefficient map near a boundary point to a Euclidean open neighborhood. Its normal component vanishes on the relevant open piece of the boundary hyperplane. An integral curve of the restricted tangent field in that hyperplane is therefore also a curve for the extended ODE. By Euclidean uniqueness a boundary initial state stays in the hyperplane, and a nearby interior solution cannot cross it: at a first hit it would agree, also for times before the hit, with the boundary solution. After shrinking the local time interval the smooth Euclidean local map restricts to a two-sided local map in the half-space. These maps are intrinsic by uniqueness; the maximal curve construction and finite-chart composition in \cref{p104:M06.03} apply to relative charts. They preserve both the boundary and the interior.

For the homotopy equivalence choose a smooth \(0\leq\lambda\leq 1/4\) equal to \(1/4\) near \(\partial M\) with closed support in the collar image, by the closed-subset cutoff of \cref{p104:M01.06}. Set
\[
 r_s(c(q,t))=c\bigl(q,t+s(1-t)\lambda(c(q,t))\bigr),\qquad 0\leq s\leq1,
\]
and use the identity off the collar. These formulas glue smoothly because the support is closed and contained in the collar image. They preserve interior points for all \(s\), and \(r_1\) maps the whole manifold into the interior. Thus \(r_1\) and the inclusion are homotopy inverses.

For a continuous map \(F\) into a manifold with boundary, first move it by \(r_s\) into the interior, then use \cref{p104:M04.05} in that boundaryless target. For a homotopy \(H\) with smooth endpoints \(F_0,F_1\), postcompose by \(r_1\) and extend the homotopy in time to \(\R\) by the endpoint maps outside \([0,1]\). This is continuous and restricts to smoothly extendable data on the closed set of the two endpoint times. The domain \(\R\times M\) has at most ordinary boundary. The relative approximation theorem gives a smooth homotopy in the interior between \(r_1F_0\) and \(r_1F_1\). Join it to the two collar homotopies, reparametrizing the time of each part by a smooth function constant near its endpoints. The joins are smooth and the original endpoints are restored. Inward-pointing fields used to construct a collar need not have a two-sided flow through the boundary; the tangency argument above is what supplies two-sided times.
\end{proof}

\section{Lie differentiation and commuting flows}

\begin{proposition}[Brackets from the flow]\label{p104:M06.07}
On a boundaryless manifold, for a vector field \(Y\) and the local flow of \(X\),
\[
 \left.\frac{\dd}{\dd t}\right|_{t=0}\varphi_t^*Y=[X,Y],
 \qquad (\varphi_t^*Y)(p)=\dd(\varphi_{-t})_{\varphi_t(p)}\,Y(\varphi_t(p)).
\]
The same derivative at other times is \(\varphi_t^*[X,Y]\). Both derivative formulas also hold for boundary-tangent \(X\) on a manifold with boundary. On a boundaryless manifold, if \([X,Y]=0\), their flows commute locally, and they commute for all times if both fields are complete. Linearly independent commuting fields are simultaneous coordinate fields near any point of independence in this setting. Compare \cite[pp.\ 227--236]{lee} and \cite[\S20.2]{tu}, including the inverse-flow base point in \cite{tuerrata}.
\end{proposition}
\begin{proof}
In coordinates \(D\varphi_t(p)=I+tDX(p)+o(t)\), by differentiating the smooth flow at zero. The factor \(\dd(\varphi_{-t})_{\varphi_t(p)}\) is its matrix inverse, with expansion \(I-tDX(p)+o(t)\). The second factor expands as \(Y(p)+tDY(p)X(p)+o(t)\). Their product gives \(DY(p)X(p)-DX(p)Y(p)\), the coordinate bracket. At time \(t\), the flow law and functoriality of pullbacks reduce the increment derivative to the same zero-time computation. For boundary-tangent \(X\), \cref{p104:M06.06} provides a smooth two-sided flow. The derivative identities extend from the dense interior by continuity in each boundary chart.

In the boundaryless setting, if the bracket is zero, \(\varphi_t^*Y=Y\) along each interval on which it is defined by integrating the zero derivative. Thus \(\varphi_t\) carries \(Y\)-curves to \(Y\)-curves along segments in its domain, and local uniqueness gives commutation for sufficiently small times near each point. For complete fields the whole segment is available for every \(t\). Conversely local commutation implies the bracket vanishes by differentiating the invariance of \(Y\). For independent commuting fields \(X_1,\ldots,X_k\), choose a transverse coordinate slice \(S\) of dimension \(n-k\). Apply the local flows successively with times \(t_1,\ldots,t_k\) to \(q\in S\). Its differential at the base point has columns the independent \(X_i\) and the slice directions, so the map is a local diffeomorphism. Commutation and the flow law permit a small increment of each time to be applied last. Differentiating shows that the corresponding coordinate field pushes forward to \(X_i\).
\end{proof}

\section{Time-dependent fields}

\begin{proposition}[Evolution maps]\label{p104:M06.08}
For a smooth time-dependent field \(X_t\) on an open subset of \(J\times M\), with \(J\) an open interval and \(M\) boundaryless, there is a unique maximal evolution \(\Psi(t,s,p)\) on its open domain. It solves the equation from \(\Psi(s,s,p)=p\), obeys \(\Psi(t,u,\Psi(u,s,p))=\Psi(t,s,p)\) on translated maximal intervals, and at fixed \(s,t\) gives inverse smooth time maps in the two directions. Compare \cite[Theorem 9.48]{lee}.
\end{proposition}
\begin{proof}
On the open subset of \(J\times M\) take the autonomous field \((1,X_t(p))\). Its maximal flow from \((s,p)\) at time \(\tau\) has time coordinate \(s+\tau\). Its state coordinate for \(\tau=t-s\) defines \(\Psi(t,s,p)\). The statements follow from \cref{p104:M06.03}, because the change \((t,s,p)\mapsto(t-s,s,p)\) is smooth and the translated time-domain equality becomes the evolution-domain equality. This construction does not assert a one-parameter action on \(M\) for a nonautonomous field.

For a left-invariant field on a Lie group, left translation carries the local solution through the identity to a solution through every other point, with the same positive and negative time widths. If a maximal interval had a finite endpoint, restarting within that uniform width would extend it. Thus the field is complete. Left invariance and the flow law give \(\gamma_e(s+t)=\gamma_e(s)\gamma_e(t)\), so its identity trajectory is a one-parameter subgroup.
\end{proof}
\chapter{Differential forms and Cartan calculus}

The exterior derivative is defined without a metric or a connection. We first fix the pointwise exterior algebra, then prove that the coordinate construction of the derivative commutes with changes of variables, including at a boundary. Throughout, local formulas are valid where the indicated charts and flow maps are defined.

\section{Exterior powers and pullback}

\begin{proposition}[The exterior algebra and its smooth sections]
\label{p104:M07.01}
Let \(V\) be a real vector space of dimension \(n\). Alternating \(k\)-linear covariant tensors on \(V\) form \(\Lambda^kV^*\). For a basis \(e^1,\ldots,e^n\) of \(V^*\), the elements
\[
 e^{i_1}\wedge\cdots\wedge e^{i_k},\qquad i_1<\cdots<i_k,
\]
form a basis. With the determinant normalization, the wedge product is associative and satisfies \(\alpha\wedge\beta=(-1)^{k\ell}\beta\wedge\alpha\) in degrees \(k,\ell\). The zeroth exterior power is \(\R\), and the exterior powers above degree \(n\) vanish. On an \(n\)-manifold \(M\), a differential \(k\)-form is a smooth section of \(\Lambda^kT^*M\). Its coefficients in the increasing coordinate wedges determine and detect its smoothness.
\end{proposition}

\begin{proof}
If \(v_1,\ldots,v_k\) are expanded in the basis dual to the \(e^i\), multilinearity and alternation collect the terms with each increasing index set into its determinant. Thus evaluation on increasing basis vectors gives exactly the coordinates of an alternating tensor. On these wedges the product of two disjoint index sets is the sorted union with the sign of the sorting permutation. A repeated index gives zero. In the basis-free shuffle formula, sum over the increasing choices \(i_1<\cdots<i_k\) of arguments for the first factor and use the increasing complement for the second factor. The sign is that of the permutation taking the original list to the concatenation of these two lists. In particular the formula agrees with determinant evaluation on products of one-forms. Sorting three sets gives the same total sign with either parenthesization, and interchanging two blocks makes \(k\ell\) transpositions. This proves the two identities by multilinearity, including degrees zero and greater than \(n\).

For a smooth frame change in \(T^*M\), the transition on exterior powers has entries that are minors of the frame change matrix. The exterior power is therefore a smooth bundle by \cref{p104:M05.04}. In a local coframe its sections have a unique coefficient expansion. The pointwise wedge is smooth on sections since the coefficient of every increasing wedge is a finite sum of products of smooth functions.
\end{proof}

This normalization agrees with \cite[Chapter 14, especially Proposition 14.11]{lee} and \cite[Sections 3 and 18]{tu}.

\begin{proposition}[Pullback and the tautological one-form]
\label{p104:M07.02}
For a smooth map \(F:M\to N\), define
\[
 (F^*\omega)_p(v_1,\ldots,v_k)
 =\omega_{F(p)}(\dd F_p v_1,\ldots,\dd F_p v_k).
\]
In degree zero this is composition with \(F\). The resulting form is smooth, pullback is contravariantly functorial, and \(F^*(\alpha\wedge\beta)=F^*\alpha\wedge F^*\beta\). Restriction to a submanifold is pullback by the inclusion. If \(M,N\) have the same dimension, the top-form coordinate multiplier is the signed Jacobian determinant. On \(T^*M\), with projection \(\pi\), there is a canonical one-form \(\vartheta\) given by \(\vartheta_\xi(V)=\xi(\dd\pi_\xi V)\).
\end{proposition}

\begin{proof}
In source coordinates \(x\) and target coordinates \(y\), \(F^*(\dd y^i)=\dd F^i\), where \(F^i=y^i\circ F\). The coefficients of the general pullback are finite sums of \(\omega_I\circ F\) times minors of \((\partial F^i/\partial x^j)\). This proves the smoothness and top-degree assertions. The chain rule proves functoriality. Substitution in the shuffle formula gives compatibility with wedge. The inclusion uses exactly the tangent inclusion of its source. Finally, in the induced cotangent coordinates \((x^i,p_i)\), where \(\xi=\sum_i p_i\dd x^i\), the projection formula is \(\vartheta=\sum_i p_i\dd x^i\). This proves smoothness of the intrinsically defined form.
\end{proof}

See \cite[Lemma 14.16, Proposition 14.20, and Chapter 22]{lee} and \cite[Sections 17.5, 17.6 and 18.5]{tu}.

\section{Exterior differentiation}

\begin{theorem}[Exterior differentiation, naturality, and support]
\label{p104:M07.03}
There is a unique family of real-linear operators \(\dd:\Omega^k(M)\to\Omega^{k+1}(M)\) that agrees with the differential on functions, satisfies \(\dd^2=0\), and satisfies
\[
 \dd(\alpha\wedge\beta)=\dd\alpha\wedge\beta+(-1)^k\alpha\wedge\dd\beta
 \quad(\alpha\in\Omega^k(M)).
\]
It is local and obeys \(F^*\dd\omega=\dd F^*\omega\) for every smooth map \(F\). In particular \(\supp(\dd\omega)\subseteq\supp(\omega)\). The coordinate formula, locality, and naturality also hold for \(C^1\) forms, with continuous derivative, and the same formula applies in smoothly compatible half-space and orthant charts.
\end{theorem}

\begin{proof}
On a coordinate domain the candidate is
\begin{equation}\label{p104:eq-coordinate-exterior}
 \dd\Big(\sum_I a_I\,\dd x^{i_1}\wedge\cdots\wedge\dd x^{i_k}\Big)
 =\sum_{I,j}\partial_j a_I\,\dd x^j\wedge\dd x^{i_1}\wedge\cdots\wedge\dd x^{i_k}.
\end{equation}
Differentiating a coefficient product and moving the new one-form past the first \(k\) slots gives the displayed Leibniz rule. The double sum for \(\dd^2\) is zero because the second derivatives are symmetric in their indices and the corresponding wedge is antisymmetric. Both identities hold first in Euclidean open sets. On a relatively open part of a half-space or orthant, local coefficient extensions give the same derivatives on the ordinary interior. That interior is dense, so the derivatives agree everywhere by continuity. This establishes extension independence, also for \(C^1\) coefficients. On a zero-dimensional model there are no positive-degree coordinates and the derivative is zero.

For the naturality calculation, use both coordinate systems on a sufficiently small source neighborhood. Each target term \(a_I\,\dd y^I\) pulls back to
\[
 (a_I\circ F)\,\dd F^{i_1}\wedge\cdots\wedge\dd F^{i_k}.
\]
Differentiating it yields the differential of \(a_I\circ F\) wedged with the same one-forms, because each \(\dd(\dd F^i)\) vanishes. The chain rule identifies that first differential with \(F^*\dd a_I\). If the target point is on a boundary or corner, use extensions of the target coefficient and the coordinate functions of \(F\) for this calculation. The extended coordinate map need only be defined near the source point and have its image in the open neighborhood on which the target coefficient was extended. It is required to map into the actual target model only on the original source. The equality thus restricts to the desired equation. The calculation uses only first derivatives of \(a_I\). In particular it proves compatibility of \eqref{p104:eq-coordinate-exterior} on chart overlaps before the global operator is defined, and proves naturality of that global operator.

Formula \eqref{p104:eq-coordinate-exterior} vanishes on an open set on which the coefficients vanish. This proves locality and the support assertion. To establish uniqueness, first show that any operator \(D\) with the stated axioms is local. If a global form \(\omega\) vanishes on an open neighborhood of \(p\), multiply it by a bump equal to one near \(p\) and supported in that neighborhood. Evaluating the Leibniz rule for this zero product at \(p\) gives \(D\omega(p)=0\). To evaluate \(D\) on a coordinate expansion, extend coefficients and coordinate functions near \(p\) using a bump supported in the chart. The local values of their derivatives are unchanged. Applying the Leibniz rule to these globally defined functions and their differentials, using \(D(\dd x^i)=0\) locally, forces \eqref{p104:eq-coordinate-exterior}.
\end{proof}

For the local-to-global argument compare \cite[Theorem 14.24 and Proposition 14.26]{lee} and \cite[Section 19]{tu}. In particular a \(C^1\) form restricts to a \(C^1\) form on a smooth face and its derivative restricts by the same identity. For any continuous form, \(\supp(F^*\omega)\subseteq F^{-1}(\supp\omega)\), because \(F\) is continuous and outside this preimage the form vanishes locally. If \(F\) is proper, this implies compactness of pullback support for compactly supported \(\omega\). For a closed embedded face the inclusion is proper.

For the tautological form from \cref{p104:M07.02}, the formula gives \(\dd\vartheta=\sum_i\dd p_i\wedge\dd x^i\) in cotangent coordinates. Thus this coordinate expression also defines a global two-form.

We will use the invariant evaluation formula
\begin{align}
 (\dd\omega)(X_0,\ldots,X_k)
 &=\sum_i(-1)^i X_i\big(\omega(X_0,\ldots,\widehat X_i,\ldots,X_k)\big)
 \notag\\
 &\quad+\sum_{i<j}(-1)^{i+j}
 \omega([X_i,X_j],X_0,\ldots,\widehat X_i,\ldots,\widehat X_j,\ldots,X_k).
 \label{p104:eq-invariant-exterior}
\end{align}
To verify it, the right side is alternating by reindexing. If \(X_0\) is multiplied by a function \(h\), the terms involving \(X_j(h)\), \(j>0\), in the first sum have coefficient \((-1)^j\). In the second sum, \([hX_0,X_j]=h[X_0,X_j]-X_j(h)X_0\) contributes the opposite coefficients. Hence the right side is linear over smooth functions in each slot, and it suffices to test it on the commuting coordinate fields. There it gives \eqref{p104:eq-coordinate-exterior}. See \cite[Propositions 14.29 and 14.32]{lee}.

\section{Contraction and Lie differentiation}

\begin{theorem}[Cartan identities]
\label{p104:M07.04}
Let \(X,Y\) be smooth vector fields on a boundaryless manifold, or tangent to the boundary in the boundary case. For forms, first-slot contraction \(\iota_X\) has degree \(-1\) and vanishes on functions. The local-flow definition of \(\mathcal L_X\) gives
\begin{align*}
 \mathcal L_X&=\dd\iota_X+\iota_X\dd, &
 [\mathcal L_X,\dd]&=0, &
 [\mathcal L_X,\iota_Y]&=\iota_{[X,Y]},\\
 [\mathcal L_X,\mathcal L_Y]&=\mathcal L_{[X,Y]}, &
 \iota_X\iota_Y+\iota_Y\iota_X&=0 .
\end{align*}
The brackets on degree-zero operators and their displayed partners are ordinary commutators. For a field on a manifold with boundary that is not tangent there, the identities remain local identities on the ordinary interior and extend as coordinate identities to the boundary.
\end{theorem}

\begin{proof}
Splitting the shuffle sum according to which factor receives the first argument gives
\[
 \iota_X(\alpha\wedge\beta)
 =(\iota_X\alpha)\wedge\beta+(-1)^k\alpha\wedge\iota_X\beta .
\]
The second contraction anticommutes by alternation. Pullbacks preserve wedges, so differentiating the local-flow pullbacks gives a degree-zero derivation. In a fixed chart the coefficients are smooth in \((t,p)\), so spatial and time derivatives commute. We may therefore differentiate the commutation of pullback with \(\dd\). On functions \(\mathcal L_X f=Xf\). The operator \(\dd\iota_X+\iota_X\dd\) is a degree-zero derivation by the two graded Leibniz identities, with the mixed terms cancelling. It gives \(Xf\) on functions and \(\dd(Xf)\) on \(\dd f\). So does \(\mathcal L_X\). Since functions and coordinate differentials generate the forms locally, the two operators coincide.

For the commutator with contraction, pullback by a local diffeomorphism \(\psi\) obeys
\[
 \psi^*(\iota_Y\omega)=\iota_{\psi^*Y}(\psi^*\omega),\qquad
 (\psi^*Y)_p=\dd(\psi^{-1})_{\psi(p)}Y_{\psi(p)}.
\]
The evaluation point of the inverse derivative is essential. This equality follows by substituting all tangent vectors. In local coordinates the derivative of \(\varphi_t^*Y\) at zero is \(DY\,X-DX\,Y=[X,Y]\). Differentiation proves the third identity. Using the Cartan formula and \([\mathcal L_X,\dd]=0\), the commutator of the two Lie derivatives is
\(\dd[\mathcal L_X,\iota_Y]+[\mathcal L_X,\iota_Y]\dd=\mathcal L_{[X,Y]}\).
All computations can be made on a fixed smaller flow neighborhood of a point for a short time. At a boundary, tangent fields have such two-sided local flows by \cref{p104:M06.06}. For other fields the invariant formula \(\dd\iota_X+\iota_X\dd\) and the component formulas on the dense interior give the asserted extension without requiring a boundary-preserving flow.
\end{proof}

In particular the formula
\[
 (\mathcal L_X\omega)(Y_1,\ldots,Y_k)
 =X\big(\omega(Y_1,\ldots,Y_k)\big)
   -\sum_i\omega(Y_1,\ldots,[X,Y_i],\ldots,Y_k)
\]
follows by writing the scalar pullback as the evaluation of the pulled-back form on the pulled-back fields and differentiating. This agrees with \cite[Propositions 14.33 and 14.34, Theorem 14.35 and Corollary 14.36]{lee}, \cite[Section 20]{tu}, and \cite[pp.~225, 232]{tuerrata}.

\section{One-forms along paths}

\begin{theorem}[Path integration and exact one-forms]
\label{p104:M07.05}
For a smooth one-form \(\alpha\), integration over a piecewise smooth path \(\gamma:[a,b]\to M\) is
\(\int_\gamma\alpha=\int_a^b\alpha_{\gamma(t)}(\dot\gamma(t))\,\dd t\), summed over its smooth pieces. It is additive under concatenation, changes sign on reversal, and is preserved by piecewise smooth nondecreasing reparametrizations onto the original interval. For a smooth function \(u\), \(\int_\gamma\dd u=u(\gamma(b))-u(\gamma(a))\). On any manifold, \(\alpha\) is exact if and only if its integral on every piecewise smooth closed path is zero, or equivalently path integrals are independent of paths with fixed endpoints. A primitive is unique up to locally constant functions.
\end{theorem}

\begin{proof}
On each smooth piece the integral is the integral of the coefficient of \(\gamma^*\alpha\), so the chain rule and one-dimensional substitution give subdivision and reparametrization invariance. Splitting intervals gives concatenation, and reversing them changes the sign. For \(\dd u\), use the derivative of \(u\circ\gamma\) and sum the fundamental theorem of calculus over the pieces.

The closed-path condition gives path independence by following one path and then the reverse of another. Any two points in a connected component can be joined by a piecewise smooth path. Indeed the set reachable from one point by such paths and its complement are both open in that component, by the existence of small convex coordinate neighborhoods, including half-balls. Choose a base point on each component and define \(u(p)\) to be the integral from that point to \(p\). In a convex relative coordinate ball about \(x\) with coefficients \(\alpha_i\), path independence gives
\[
 u(y)-u(x)
 =\int_0^1\sum_i \alpha_i(x+t(y-x))(y^i-x^i)\,\dd t .
\]
The right side is smooth in \(y\) using local extensions in a boundary chart and differentiation under the integral. For every interior \(x\), the derivative at \(y=x\) is \(\sum_i\alpha_i(x)\dd x^i\). Equality extends to the boundary by continuity. Thus \(\dd u=\alpha\). If two primitives differ, their difference is constant on piecewise smooth paths by the first part, hence on every component. The empty manifold is covered by the same statement with no components to choose.
\end{proof}

A closed form may have a nonzero period. On the unit circle, the restriction of \(x\,\dd y-y\,\dd x\) is closed for dimensional reasons and integrates to \(2\pi\) on one positive turn. This rules out a primitive by the theorem. The fixed-endpoint homotopy invariance of closed one-form integrals will follow from Stokes. See \cite[Chapter 11, pp.~287--298]{lee} and \cite[Sections 18.7 and 19.7]{tu}.

\section{Lie differentiation of general tensors}

\begin{proposition}[Tensor Lie derivatives]
\label{p104:M07.06}
Pullback along the local diffeomorphisms of a field's flow defines Lie differentiation on smooth tensors of arbitrary variance wherever that flow is defined. It agrees with \(Xf\) on functions and with \([X,Y]\) on vector fields, acts as a derivation for tensor products, and commutes with contraction and permutations of slots. In particular, for a covariant tensor \(A\),
\[
 (\mathcal L_X A)(Y_1,\ldots,Y_s)
 =X\big(A(Y_1,\ldots,Y_s)\big)
  -\sum_{i=1}^s A(Y_1,\ldots,[X,Y_i],\ldots,Y_s).
\]
These formulas are independent of a connection and extend to a boundary. Compare \cite[Chs.~9 and 12]{lee} and \cite[\S20]{tu}.
\end{proposition}

\begin{proof}
A diffeomorphism transports tangent vectors by its inverse derivative at the image point and transports covectors by the dual of its forward derivative. The resulting pointwise maps preserve vector--covector evaluation, tensor products, contractions, and permutations. Differentiation gives the stated algebraic identities. For vector fields the derivative is \(DY\,X-DX\,Y\) as in the proof of \cref{p104:M07.04}. Evaluating the pullback on transported fields and differentiating proves the displayed formula. In local bases, mixed-variance tensors are finite sums of tensor products of coordinate fields and covectors with smooth coefficients. The formulas for those factors and the product rule then express every derivative smoothly in the coefficients and their first derivatives. For a boundary-tangent field this matches the actual flow calculation up to the boundary. For a general field, agreement on the interior shows uniqueness of the smooth local extension to a boundary chart.
\end{proof}
\chapter{Distributions, foliations, and smooth quotients}

\section{Involutivity and annihilating forms}

\begin{definition}
A smooth distribution of rank \(k\) on a boundaryless manifold \(M\) is a smooth rank-\(k\) vector subbundle \(D\subseteq TM\). It is involutive if the bracket of any two local smooth sections of \(D\) is a section of \(D\). An immersed integral submanifold has dimension \(k\) and its tangent spaces map onto the corresponding fibers of \(D\). The distribution is integrable if an integral submanifold passes through every point.
\end{definition}

\begin{proposition}[Bracket and annihilator tests]
\label{p104:M08.01}
An integrable distribution is involutive. Involutivity can be checked on a local frame of \(D\). Write \(D^\circ\subseteq T^*M\) for the annihilator subbundle. The distribution is involutive if and only if the exterior derivative of each local section \(\theta\) of \(D^\circ\) vanishes on pairs of vectors of \(D\). Equivalently, the graded ideal of forms that vanish when restricted to \(D\) is stable under \(\dd\).
\end{proposition}

\begin{proof}
Near any point on an integral submanifold, its inclusion is an embedding after restricting the domain. A smooth field with values in \(D\) is tangent to that embedded patch. The bracket of two such fields is related to the bracket of their restrictions under the inclusion, by naturality of the bracket for related fields. Its value is therefore in \(D\) at that point. A local frame suffices because
\([fX,gY]=fg[X,Y]+fX(g)Y-gY(f)X\).

The annihilator is a subbundle by the dual frame construction. For local sections \(X,Y\) of \(D\) and \(\theta\) of \(D^\circ\), \eqref{p104:eq-invariant-exterior} gives
\[
 \dd\theta(X,Y)=-\theta([X,Y]).
\]
Every covector in the annihilator extends to a local section of the annihilator using a local frame, so simultaneous vanishing of these evaluations is equivalent to \([X,Y]\in D\). Extend a local frame of \(D\) to one of \(TM\). In its dual coframe a \(j\)-form restricts to zero on \(D\) exactly when its expansion has at least one annihilator covector in every nonzero term. Thus the local ideal is generated by the annihilator one-forms. Involutivity puts the derivative of each generator in the ideal. The Leibniz rule then proves preservation of the ideal. Conversely, preservation of the global ideal gives the local test by multiplying a local annihilator section by a bump equal to one near the test point and extending it by zero. Its derivative at the test point equals that of the original local section.
\end{proof}

See \cite[Proposition 19.3, Lemma 19.4, Theorem 19.7 and Proposition 19.11]{lee}.

\section{The local Frobenius theorem}

\begin{theorem}[Adapted coordinates for an involutive distribution]
\label{p104:M08.02}
For a smooth rank-\(k\) distribution on an \(n\)-manifold without boundary, involutivity is equivalent to the existence, near every point, of coordinates \((u^1,\ldots,u^k,z^1,\ldots,z^{n-k})\) in which \(D\) is spanned by the \(u\)-coordinate fields. In these coordinates the slices \(z=\mathrm{constant}\) are integral submanifolds.
\end{theorem}

\begin{proof}
Take a chart centered at \(p\) and linearly change its coordinates so the projection \(\pi\) to the first \(k\) coordinates has differential an isomorphism from \(D_p\) to \(\R^k\). This remains true on a smaller chart because the determinant of the projected local frame is nonzero at \(p\). There are unique smooth sections \(Y_i\) of \(D\) whose projected vectors are the constant coordinate vectors of \(\R^k\). The \(Y_i\) are \(\pi\)-related to commuting vector fields. Thus \([Y_i,Y_j]\) lies in \(\ker\dd\pi\). It also belongs to \(D\) by involutivity. Injectivity of \(\dd\pi\) on \(D\) forces all these brackets to vanish.

Let \(S\) be the coordinate slice on which the first \(k\) coordinates vanish. On a small product neighborhood of \((0,p)\) in \(\R^k\times S\), define \(\Psi(u,q)\) by flowing from \(q\) for time \(u^i\) in \(Y_i\), in succession. A common neighborhood for this finite composition exists by openness of the local flow domains. After shrinking it, the commutation of the local flows from \cref{p104:M06.07} allows any one flow to be placed last. Consequently \(\partial\Psi/\partial u^i=Y_i\circ\Psi\). At \((0,p)\), the \(u\)-derivatives are a basis of \(D_p\) and the derivatives in \(q\) are the complementary vectors in \(T_pS\). The inverse function theorem gives the desired chart. The converse follows either from the frame test or from the slices themselves. For \(k=0\) the coordinates have only transverse variables and the slices are points. For \(k=n\) the transverse factor is a point and the whole chart is integral.
\end{proof}

For this local construction see \cite[Theorem 19.12 and Corollary 19.13]{lee}.

\section{Leaves and their topology}

\begin{theorem}[Global leaves of a regular distribution]
\label{p104:M08.03}
An involutive smooth constant-rank distribution on a boundaryless manifold partitions it into unique maximal connected immersed integral submanifolds. These leaves admit compatible local plaques in the adapted charts and, with their own manifold topology, are Hausdorff and second countable. A leaf can have a topology strictly finer than its subspace topology.
\end{theorem}

\begin{proof}
Choose a countable cover by adapted charts whose images are products of open balls. Within one chart call each slice with fixed transverse coordinate a plaque. Each plaque is an embedded connected integral patch. Declare two points equivalent when a finite chain of intersecting plaques joins them. A plaque that meets a second adapted chart locally has constant transverse coordinate in that chart, because the derivative of the second transverse coordinate vanishes on \(D\). In particular, around each intersection point two plaques through it agree locally as subsets of \(M\). To see both inclusions, apply this argument first to a small connected part of the first plaque lying in the second chart, then to a small connected part of the second lying in the first chart. Projection onto the appropriate longitudinal coordinates shows that the overlap is open in both plaques and that the restricted transition maps are smooth.

Equip an equivalence class \(L\) with the topology whose basis consists of open subsets of its plaques. The overlap observation supplies the basis property and compatible smooth transitions. Ambient open sets restrict to open sets in \(L\), so the Hausdorff property follows from that of \(M\). Each plaque meets at most countably many plaques of the fixed cover. Indeed, its intersection with any one chart is open in the plaque and has at most countably many components, since components of a Euclidean open set are open and a countable basis detects them. Each component has constant transverse coordinate in the intersecting chart, so lies in just one of its plaques. The set of plaques reachable from a fixed plaque by chains of any prescribed finite length is therefore countable by induction. Taking the countable union over lengths shows that the leaf has a countable plaque cover, and the Euclidean bases on these plaques give second countability. This argument does not assume a leaf meets an entire chart in only one plaque.

The inclusion \(L\to M\) is a smooth injective immersion by the plaque charts. The class \(L\) is connected because it is the union of connected plaque chains meeting a fixed plaque. If \(f:N\to M\) is any connected integral immersion, then in each adapted chart the transverse coordinate of \(f\) is locally constant on its inverse image, since its differential is zero. Thus the image near any point of \(N\) lies in one plaque. The inverse images of different classes are disjoint open sets and \(N\) is connected, so the image lies in a single \(L\). The local factorization through plaques is smooth and has invertible differential, since the dimensions agree. This proves maximality and the uniqueness of the connected leaf structure.
\end{proof}

For an explicit example, let \(D\) on \(\R^2/\mathbb Z^2\) be the span of the constant field of slope \(\alpha\notin\mathbb Q\). The leaf through the origin is parametrized injectively by \(t\mapsto(t,\alpha t)\bmod\mathbb Z^2\). Its image is a whole leaf because the field's orbits are open in each connected one-dimensional leaf and partition it. The parametrization is a local diffeomorphism onto the leaf and hence a diffeomorphism for the leaf topology. The fractional parts of \(j\alpha\) for integer \(j\) are dense in the circle. One proof uses the pigeonhole principle on the first \(N+1\) multiples to obtain a nonzero step of size less than \(1/N\) modulo integers, and then uses integer multiples of that step to approximate an arbitrary point on the circle. Fixing the first coordinate modulo one shows that the leaf is dense in the torus. There are arbitrarily large positive integers \(j\) for which \(j\alpha\) modulo one is arbitrarily close to zero, since the small nonzero steps from increasing \(N\) cannot all come from a fixed finite set of integers. Thus ambient convergence to the origin along the leaf does not imply convergence of its real parameter to zero. This separates the two topologies. See \cite[pp.~501--505, especially Theorem 19.21]{lee}.

\section{Free proper actions and their orbit manifolds}

\begin{theorem}[Smooth quotient of a free proper action]
\label{p104:M08.04}
Let a finite-dimensional Lie group \(G\) act smoothly, freely, and properly on a nonempty boundaryless \(n\)-manifold \(M\). Properness means that
\[
 G\times M\longrightarrow M\times M,\qquad (g,p)\longmapsto (p,gp)
\]
is proper. The orbit space \(M/G\) has a unique smooth manifold structure of dimension \(n-\dim G\) for which the quotient map is a smooth submersion. A map out of \(M/G\) to a smooth manifold is smooth precisely when its composite with the quotient projection is smooth. If \(G\) is discrete the projection is a smooth covering.
\end{theorem}

\begin{proof}
Write \(q:M\to M/G\). This map is open because the inverse image of \(q(U)\) is the union of the translates of \(U\). The orbit relation is closed in \(M\times M\). Indeed, for any convergent sequence \((p_j,g_jp_j)\to (p,p')\), properness over the compact set consisting of this sequence and its limit gives a convergent subsequence of \((g_j,p_j)\). Its limit \((g,p)\) satisfies \(gp=p'\). Manifold first countability suffices for the resulting sequential closedness argument. If two points represent different orbits, choose a product of open neighborhoods of them disjoint from the orbit relation. Their images are disjoint open neighborhoods of the two orbits. Also the images of a countable basis of \(M\) form a basis of \(M/G\), because \(q\) is open and onto. This proves the required topological properties.

The differential at the identity of the orbit map \(a_p(g)=gp\) is injective. If its value on \(v\in T_eG\) is zero, take the short integral curve \(g(t)\) through \(e\) of the left-invariant vector field determined by \(v\). Write \(\lambda_g(x)=gx\). The derivative of the orbit curve is
\[
 \frac{\dd}{\dd t}\big(g(t)p\big)
 =(\dd\lambda_{g(t)})_p\,(\dd a_p)_e v=0.
\]
So \(g(t)p=p\) and freeness implies \(g(t)=e\) throughout this interval. Differentiating at zero gives \(v=0\).

Take a local embedded slice \(S_0\) through \(p\) with \(T_pS_0\) complementary to the image of \((\dd a_p)_e\), by selecting a coordinate plane with that tangent space. The action map \(\Phi(g,s)=gs\) has invertible derivative at \((e,p)\) on \(G\times S_0\). The inverse function theorem allows neighborhoods \(W\) of \(e\) and \(S_1\) of \(p\) in \(S_0\) on whose product \(\Phi\) is a diffeomorphism onto an open set. Properness and freeness now give a neighborhood \(U\) of \(p\) such that if \(x,y\in U\) and \(gx=y\), then \(g\in W\). Otherwise take \(x_j,y_j\to p\) in a decreasing neighborhood basis and \(g_j\notin W\) with \(g_jx_j=y_j\). Properness gives a subsequential limit \((g,p)\), with \(g\notin W\) and \(gp=p\). This contradicts \(g=e\) by freeness.

Shrink the slice to \(S=S_1\cap U\). If \(s,t\in S\) and \(gs=t\), then \(g\in W\) and injectivity on \(W\times S_1\) implies \((g,s)=(e,t)\). Thus each orbit meeting \(S\) does so in exactly one point, and \(\Phi:G\times S\to GS\) is bijective. Its derivative is invertible everywhere because this is true at \((e,s)\) and the action of each fixed \(g\) and left translation in \(G\) are diffeomorphisms. Hence \(GS\) is open and the bijective local diffeomorphism \(\Phi\) is a diffeomorphism.

The diffeomorphism gives a quotient chart identifying \(q(GS)\) with \(S\) and identifies \(q\) with projection to the slice. For two slices the transition sends \(s\) to the slice component of \(\Phi'^{-1}(s)\) in the other trivialization wherever \(s\in GS'\). This is smooth, and the same argument gives a smooth inverse. The projection is therefore a submersion. It has smooth local sections in these charts. Local sections also exist for any other structure making it a submersion, by the local submersion theorem. These sections prove the universal property for smooth maps and smoothness of the identity in both directions between any two such structures. Finally, when \(G\) is discrete, \(G\times S\) is the disjoint union of open copies of \(S\), each projecting diffeomorphically onto \(q(GS)\). This is the covering condition.
\end{proof}

For example the free translation action of \(\mathbb Z^n\) on \(\R^n\) is proper since bounded \(p\) and \(p+g\) allow only finitely many integer vectors \(g\). Its quotient is the smooth torus. The theorem also applies to the antipodal action on \(S^n\) for \(n\geq1\), giving the smooth projective-space quotient. The quotient charts of this example agree with the projective charts by the local-section characterization. See \cite[Proposition 21.5 and Theorems 21.10 and 21.13]{lee} and \cite[Section 7]{tu}.
\chapter{Orientations and boundary signs}

An orientation is a choice of positive ray in a determinant line, varying continuously with the base point. The determinant line formulation covers zero-dimensional manifolds, boundary charts, and the orientation covering with the same convention.

\section{Equivalent descriptions of orientation}

\begin{theorem}[Orientations and positive top forms]
\label{p104:M09.01}
For a smooth \(n\)-manifold \(M\) the following data are equivalent.
\begin{enumerate}[label=\textup{(\alph*)}]
 \item a continuous choice of positive ray in \(\Lambda^nT_pM\) at each \(p\)
 \item smooth local generators of the determinant line, related on overlaps by positive factors
 \item a compatible assignment of signs to local coordinate frames
 \item a nowhere vanishing smooth top form, taken up to multiplication by an everywhere positive smooth function.
\end{enumerate}
In dimension zero an orientation is a sign at each point. If \(M\) is connected and nonempty and admits an orientation, it has exactly two orientations.
\end{theorem}

\begin{proof}
A local frame trivializes the determinant line, and two frames represent the same ray exactly when their transition matrix has positive determinant. For \(n=0\), use the canonical identification of the zeroth exterior power with \(\R\) and keep the sign rather than asking the empty basis to realize both rays. More generally a chart with coordinate frame \(E=(E_1,\ldots,E_n)\) has an orientation sign \(\epsilon_U\) defined by requiring \(\epsilon_U(E_1\wedge\cdots\wedge E_n)\) to lie on the positive ray. This sign is locally constant. On an overlap it changes by the sign of the Jacobian of the transition map. Conversely this compatibility makes the local choices of ray agree. For \(n>0\) interior charts can be made positive by reflecting one coordinate when necessary. A signed-chart convention also works in a fixed half-space model, where the allowed chart reflection need not exist in dimension one.

A top form \(\omega\) is positive if its pairing with every positive determinant vector is positive. Local positive top forms come from the dual coordinate wedges with the chosen sign. Multiply them by a smooth partition of unity subordinate to these charts and extend each product by zero. Their sum is smooth by local finiteness. At each point at least one coefficient in the partition is positive, and all nonzero summands are positive on the same ray. This yields a positive smooth global top form. Conversely such a form defines the positive ray by its pairing. Two forms define the same orientation exactly when their smooth pointwise ratio is positive. Finally the pointwise comparison of two orientations is a locally constant sign. It is constant on a connected nonempty manifold. The construction applies componentwise and gives the empty manifold its unique empty choice.
\end{proof}

This description uses \cite[Chapter 15, pp.~378--388]{lee} and \cite[Section 21]{tu}. In the signed-coordinate formulation, the sign belongs to the chart convention, and the associated orientation is intrinsic.

\section{Products and the orientation double cover}

\begin{proposition}[Products, local signs, and the orientation covering]
\label{p104:M09.02}
The product orientation concatenates positive determinant vectors with the first factor before the second. A local diffeomorphism between oriented manifolds has a locally constant sign, and signs multiply under composition. Every manifold \(M\) has a canonical smooth two-sheeted covering \(\pi:\widetilde M\to M\), whose points are pairs \((p,o_p)\) with \(o_p\) an orientation of \(T_pM\). The cover is orientable. For connected nonempty \(M\), the cover has two connected components if \(M\) is orientable and is connected if \(M\) is not orientable.
\end{proposition}

\begin{proof}
The direct sum decomposition of a product tangent space identifies its determinant line with the tensor product of the factor determinant lines in the indicated order. A change of bases acts by the product of determinants. This gives a coordinate-independent product orientation and includes the cases where one dimension is zero. For a local diffeomorphism the derivative on determinant lines is a continuous family of isomorphisms, so its sign relative to the two rays is locally constant. The determinant of a composition is the product of the determinants. These arguments also apply in product corner charts.

For the cover, choose any local frame and use it to identify the two orientations over its domain with the two labels \(+\) and \(-\). On overlaps the label is multiplied by the locally constant transition sign of the two frames. These identifications give compatible charts, and their projection is a local diffeomorphism onto the original chart. Two points over distinct base points can be separated by disjoint inverse images of base neighborhoods. Two points over the same base point are separated by the two sheets in a local trivialization. A countable base cover with countable bases on both sheets gives second countability. At \((p,o_p)\), pull back \(o_p\) by \(\dd\pi\) to orient the cover. This orientation is smooth in the local sheets. The deck involution exchanging the labels reverses it.

An orientation of \(M\) defines a smooth section of the cover, and the section and its negative give the two disjoint open copies of \(M\) in the orientable case. Conversely, assume \(M\) is connected and the cover is disconnected. Connected manifolds are path connected by local path connectedness. Starting at any point in a component of the cover, lift paths to see that the component projects onto \(M\). There can therefore be at most two components. If there are two, each meets each fiber in one point and projects by a bijective local diffeomorphism. Its inverse is a smooth section of the covering and gives an orientation of \(M\).
\end{proof}

For example, \(S^2\) has its orientation induced as the boundary of the unit ball in \(\R^3\). The antipodal map is the restriction of \(-I\), which sends outward normals to outward normals and has determinant \(-1\) in \(\R^3\). Thus the antipodal map reverses the sphere orientation. If \(\mathbb{RP}^2\) had an orientation, pulling it back by the two-sheeted local diffeomorphism \(S^2\to\mathbb{RP}^2\) would give a deck-invariant orientation of \(S^2\). Every orientation on the connected \(S^2\) is one of its two opposite orientations, both reversed by the antipodal map. This proves nonorientability of \(\mathbb{RP}^2\). See \cite[Proposition 15.7 and Theorem 15.41]{lee} and \cite[Section 21]{tu}.

\section{The outward-first convention}

\begin{theorem}[Induced orientation on a boundary]
\label{p104:M09.03}
Let \(M\) be an oriented smooth \(n\)-manifold with boundary and \(n\geq1\). There is a canonical orientation on \(\partial M\). At \(p\in\partial M\), a determinant vector \(\eta\) for \(T_p\partial M\) is positive if \(v\wedge\eta\) is positive for \(T_pM\), where \(v\in T_pM\) is any outward-pointing vector. For a half-space chart \(x_n\geq0\) with orientation sign \(\epsilon_U\), the ordered tangential coordinate directions carry boundary sign \(\epsilon_U(-1)^n\). The positively oriented interval has sign \(-1\) at its left endpoint and \(+1\) at its right endpoint.
\end{theorem}

\begin{proof}
The line \(T_pM/T_p\partial M\) has a distinguished inward half-line given in boundary coordinates by increasing \(x_n\). To see its independence of chart, consider the last coordinate \(y_n\) of a transition. It vanishes on the boundary and is nonnegative on the model side, so its tangential derivatives vanish and its inward normal derivative is nonnegative. The differential of the transition is invertible and preserves the boundary tangent subspace. It follows that its induced map on the quotient line is nonzero, hence the inward normal derivative is strictly positive. Thus the opposite ray defines outwardness independently of the chart.

Replacing an outward vector by another outward vector replaces it by a positive multiple modulo \(T_p\partial M\). The tangential difference wedges to zero with \(\eta\). Thus the sign definition is independent of the outward vector. The resulting rays are smooth along the boundary in local frames. In particular the vector used in the construction is transverse to the boundary, although its coefficients can be specified only along the boundary. This differs from being a tangent field on \(\partial M\).

In the half-space chart the vector \(-\partial_{x_n}\) points outward. The wedge with the ordered tangential frame is
\[
 (-\partial_{x_n})\wedge\partial_{x_1}\wedge\cdots
 \wedge\partial_{x_{n-1}}
 =(-1)^n\partial_{x_1}\wedge\cdots\wedge\partial_{x_n}.
\]
Hence the tangential frame has the asserted boundary sign. At the left endpoint of the standard oriented interval the inward coordinate is \(x-a\) and the chart sign is \(+1\). At the right endpoint it is \(b-x\) and the chart sign is \(-1\). This gives the two endpoint signs, using the zero-dimensional orientation convention of \cref{p104:M09.01}.
\end{proof}

On a product of oriented manifolds with dimensions \(m,n\), the boundary orientation on the face \(\partial M\times \operatorname{Int} N\) is the product orientation. The face \(\operatorname{Int}M\times\partial N\) carries \((-1)^m\) times its product orientation, because the outward direction of the second factor must cross the \(m\) tangent directions of the first. This is a statement on codimension-one faces. For a product with corners, their higher-codimension intersections will be handled by localization in Stokes's theorem. See \cite[Proposition 15.24]{lee}, \cite[Sections 22.3--22.6]{tu}, and \cite[p.~254]{tuerrata}.

\section{Volume forms and densities from metrics}

\begin{proposition}[Metric volume and metric density]
\label{p104:M09.04}
A smooth positive-definite metric \(g\) on an oriented manifold \(M\) determines a unique smooth top form \(\mathrm{vol}_g\) that evaluates to one on the positive unit vector of each determinant line \(\Lambda^nT_pM\), with the norm induced by \(g\).

In signed coordinates with metric matrix \(G=(g(\partial_i,\partial_j))\),
\[
 \mathrm{vol}_g
 =\epsilon_U\sqrt{\det G}\,\dd x^1\wedge\cdots\wedge\dd x^n .
\]
On every manifold with such a metric, including a nonorientable one, the coordinate expression \(\sqrt{\det G}\,\lvert\dd x^1\cdots\dd x^n\rvert\) gives a positive smooth density.
\end{proposition}

\begin{proof}
The metric gives \(E=\partial_1\wedge\cdots\wedge\partial_n\) norm \(\sqrt{\det G}\), by the Gram determinant formula. Thus the positive unit vector is \(\epsilon_U E/\sqrt{\det G}\). The displayed form is the unique functional on the determinant line taking value one on that vector. If \(n>0\), the wedge of a positively oriented orthonormal frame is exactly this positive unit vector. For the chart change \(x=x(y)\), let \(A=(\partial x^i/\partial y^j)\). The metric matrix in the \(y\)-coordinates is \(A^{\mathsf T}G_x A\) and its square-root determinant is \(\lvert\det A\rvert\sqrt{\det G_x}\). The top-form basis changes by \(\det A\), and \(\epsilon_y=\epsilon_x\operatorname{sgn}(\det A)\). Thus the signed expressions are equal on overlaps. A density is a quantity whose coordinate coefficient changes using the absolute determinant, so the unsigned expressions also agree on overlaps, without a choice of rays. Both are smooth since the determinant of a positive-definite matrix is positive and smooth. In dimension zero \(E=1\) and the empty determinant is one. Thus the positive unit vector is the assigned point sign, the volume form takes that same signed value on \(1\), and the metric density is one.
\end{proof}

Reversing the orientation reverses \(\mathrm{vol}_g\) and gives the same density. The pointwise density of a metric is nowhere zero. Absolute values of arbitrary top forms will be treated separately in the integration chapter, where zeros affect smoothness. For the volume construction see \cite[Proposition 15.29 and pp.~427--432]{lee}.
\chapter{Integration and Stokes's theorem}

\section{Top forms and coordinate integration}

For continuous top forms with compact support, integration is obtained from
ordinary Euclidean integration. A form supported in a coordinate domain is
understood to have support compact in that domain whenever it is extended by
zero outside the coordinate image. On a half-space chart this allows support
on its true boundary, but there is an open neighborhood of every artificial
chart edge on which the extension is zero.

\begin{proposition}[The local signed substitution formula]
\label{p104:M10.01}
In an oriented \(n\)-dimensional chart, \(n\geq1\), the integral of a
compactly supported top form with coefficient \(a\) is the Euclidean
integral of \(a\), multiplied by the chart's orientation sign if its
standard coordinate orientation is negative. The result agrees on two
charts containing the support. In dimension zero the corresponding
integral is signed evaluation.
\end{proposition}
\begin{proof}
Let \(y=\psi(x)\) be the transition on the overlap. If the form has
coefficient \(b(y)\), its coefficient in \(x\) is
\(b(\psi(x))\det D\psi(x)\). Euclidean substitution uses
\(\lvert\det D\psi\rvert\); the residual sign is the ratio of the two
chart signs. If the sign is not constant on an overlap, the identity
holds separately on the open positive and negative loci. In the
half-space case the transition maps interiors to interiors, and the
coordinate boundary lies in a hyperplane of measure zero. Thus the
usual substitution theorem applies to the interiors, with unchanged
integrals. In dimension zero each point is open and has its prescribed
sign. A compact subset of a discrete space is finite, so a
compactly supported zero-form gives a finite signed sum.
\end{proof}

\begin{theorem}[Integration and independence of localization]
\label{p104:M10.02}
On an oriented manifold \(M\) with or without boundary there is a unique
linear integral of compactly supported continuous top forms compatible
with coordinate integration and finite localization. Reversing the
orientation changes its sign. A diffeomorphism \(F:M\to N\) preserving
orientation satisfies \(\int_M F^*\eta=\int_N\eta\), and if it reverses
orientation the right hand side is negated. More generally the
substitution sign is taken componentwise.
\end{theorem}
\begin{proof}
Choose a locally finite smooth partition of unity \((\rho_i)\) subordinate
to signed coordinate domains \(U_i\). Each \(\rho_i\eta\) has support
contained in \(\supp\eta\cap\supp\rho_i\), a compact subset of \(U_i\).
Only finitely many terms can meet \(\supp\eta\). Define the integral by
summing their local integrals. For a second partition \((\chi_j)\),
insert \(\sum_j\chi_j=1\) in each summand. The resulting products
\(\rho_i\chi_j\eta\) have compact support in \(U_i\cap V_j\); the preceding
coordinate formula lets us calculate each in either chart. Summing first
over the \(i\)-index gives exactly the second integral. This argument also
allows partitions summing to one on a neighborhood of the given support
and continuous localizing weights with the same subordination. A common
partition proves linearity and shows uniqueness. Changing the orientation
changes every local sign. For a diffeomorphism, pull the partition and
charts of \(N\) back to \(M\). The pullback has compact support because a
diffeomorphism is proper, and the claimed sign is precisely the sign of
the transported chart. In dimension zero the same arguments reduce to
finite sums and bijective reindexing.
\end{proof}

Neither an arbitrary smooth pullback nor restriction to an arbitrary open
submanifold preserves compact support. For a closed embedded submanifold
\(S\subset M\), restriction does preserve it: the support on \(S\) is
contained in the intersection of \(S\) with the original compact support.
This applies to the boundary inclusion. The constructions in this
section are those of \cite[Chapter 16, pp.~402--410]{lee} and
\cite[\S23.1--23.4]{tu}; in signed coordinates the determinant in the
change of variables must be kept \cite[pp.~266--267]{tuerrata}.

\section{The half-space calculation and its globalization}

\begin{lemma}[Local Stokes in all dimensions]
\label{p104:M10.03}
For \(n\geq1\), a compactly supported \(C^1\) \((n-1)\)-form on
\(\R^n\) has integral of its exterior derivative zero. On \(\HH^n\),
with the standard bulk orientation and outward-first boundary
orientation, that integral equals the integral of its restriction to
\(x_n=0\).
\end{lemma}
\begin{proof}
Write
\[
 \omega=\sum_{j=1}^n(-1)^{j-1}a_j\,
     \dd x_1\wedge\cdots\wedge\widehat{\dd x_j}\wedge\cdots\wedge\dd x_n .
\]
Then \(\dd\omega=(\sum_j\partial_j a_j)\dd x_1\wedge\cdots\wedge\dd x_n\).
The fundamental theorem on each full coordinate line, followed by Fubini,
gives zero for each full-line derivative integral. On the half-space
only \(\partial_n a_n\) remains. Its integral along the last coordinate
is \(-a_n(x',0)\). The coefficient of the restricted form in the first
\(n-1\) coordinates is \((-1)^{n-1}a_n(x',0)\). The sign of those
coordinates as an outward-first boundary frame is \((-1)^n\), so the
boundary coefficient contributes \(-a_n(x',0)\), as required.
Continuity and compact support justify Fubini even on the boundary
model. In dimension one the restriction is the value at the endpoint,
with the negative sign for the left endpoint of a positively oriented
interval.
\end{proof}

\begin{theorem}[Stokes on an abstract manifold]
\label{p104:M10.04}
Let \(M\) be an oriented smooth \(n\)-manifold with boundary, \(n\geq1\),
with outward-first boundary orientation. If \(\omega\) is a \(C^1\)
\((n-1)\)-form with compact support in \(M\), then
\[
                  \int_M \dd\omega=\int_{\partial M}j^*\omega ,
                  \qquad j:\partial M\hookrightarrow M .
\]
For a compact \(M\) the support condition is automatic, and for an
empty boundary the right side is zero.
\end{theorem}
\begin{proof}
The exterior derivative in charts is continuous, is supported in
\(\supp\omega\), and agrees across charts. For \(C^1\) forms the latter
claim follows directly from the coordinate formula and smoothness of
transitions: write a form as a sum of \(a_I\dd x_I\), differentiate the
\(C^1\) coefficients after substitution, and note that derivatives of
the coordinate differentials have vanishing skew part. The same
calculation gives naturality under restriction to a face. At a
half-space boundary the \(C^1\) formula follows from locally extended
coefficients and continuity from the interior.

Choose a subordinate partition \((\rho_i)\) as above. The products
\(\rho_i\omega\) extend by zero in their Euclidean models and are still
\(C^1\) up to the true boundary. No term meets an artificial edge in
its coordinate model. The finitely many nonzero terms satisfy
\[
   \sum_i \dd(\rho_i\omega)
     =\sum_i\bigl(\dd\rho_i\wedge\omega+\rho_i \dd\omega\bigr)
     =\dd\omega ,
       \qquad
   \sum_i j^*(\rho_i\omega)=j^*\omega .
\]
These identities hold also at support boundary points, since they
come from a locally finite smooth sum equal to one on \(M\).
The boundary is closed in \(M\), and the boundary restrictions have
compact support there. Apply the local formula to each term, with the
same orientation sign on the bulk and its induced boundary. The
coordinate definition of integration and linearity finish the proof.
\end{proof}

For example, if \(D\) is a compact smooth planar domain with the
orientation \(\dd x\wedge\dd y\), apply the theorem to
\(P\,\dd x+Q\,\dd y\). Its derivative has coefficient
\(\partial_xQ-\partial_yP\); the outer boundary is counterclockwise and
the inner boundaries have the opposite sign. In an oriented Euclidean
domain of any dimension, \(\omega=\iota_X(\dd x_1\wedge\cdots\wedge\dd
x_n)\) has derivative \((\operatorname{div}X)\dd x_1\wedge\cdots\wedge\dd
x_n\), and its boundary restriction is the outward flux. On an oriented
surface in \(\R^3\), the pullback of the derivative of the one-form
associated to a vector field is the normal component of its curl times
area. These formulas use the same outward-first convention and follow
from evaluating the indicated forms in coordinate frames. See
\cite[Theorem 16.11 and pp.~411--415, 421--425]{lee} and
\cite[\S23.5]{tu}.

\section{Densities and corner faces}

\begin{theorem}[Density integration]
\label{p104:M10.05}
A density on an \(n\)-manifold is a section of the real density line,
whose coefficient in a new frame changes by the absolute determinant.
Continuous compactly supported densities have a linear integral on any
manifold, invariant under diffeomorphisms, without an orientation.
For an equal-dimensional smooth \(F\), pullback of a continuous density
is continuous; if the density is smooth and \(F\) is a local
diffeomorphism its pullback is smooth.
\end{theorem}
\begin{proof}
At a tangent space \(V\), a density is specified by its value on a
basis, with the rule \(\mu(A v)=|\det A|\,\mu(v)\) for a basis
transformation. Extend its frame evaluation to linearly dependent
lists by zero. The multiplicative determinant identity gives the
transition law of a line bundle. The coordinate coefficient
\(b(x)|\dd x_1\cdots\dd x_n|\) transforms by
\(b(F(x))|\det DF(x)|\) under pullback. This expression remains
continuous at a rank drop. If the determinant is nonzero it has a
constant sign locally, so its absolute value is smooth. Integrate in
coordinates and use ordinary absolute-Jacobian substitution. The
common-refinement argument for top forms now has no signs. The
zero-dimensional rule is the sum of the coefficients on the
finitely many support points.
\end{proof}
A smooth Riemannian metric gives the positive density
\(\sqrt{\det(g_{ij})}|\dd x_1\cdots\dd x_n|\). Its intrinsic description
is the volume of the parallelepiped given by a frame, which obeys the
same transformation law. In contrast, \(x\mapsto x^2\) pulls back
\(|\dd y|\) to \(2|x||\dd x|\). Likewise the pointwise absolute value
of a smooth top form can fail to be smooth at a zero. These
regularity distinctions are important in \cite[Propositions
16.35--16.42]{lee} with the correction at p.~430 of \cite{leeerrata}.

For corner models we use relatively open subsets of
\(\R_+^k\times\R^{n-k}\) with smoothly extendible coordinate transitions
and inverses. At a point \(p\), let \(I(p)\) be the set of the \(k\)
constrained coordinates vanishing there. The velocity cone of curves
with parameter \(t\geq0\) at \(p\) is exactly
\(\{v:v_j\geq0\text{ for }j\in I(p)\}\), by differentiating
constraints in one direction and by straight rays in the other.
The derivative of a transition sends this cone isomorphically to the
new cone. Its maximal linear subspace has codimension \(|I(p)|\);
this proves that the number of active faces is invariant. The facets
of the cone are its codimension-one faces, and the same isomorphism
permutes them. In a neighborhood of \(p\), each facet germ corresponds
to a local boundary hypersurface through \(p\). For example, the
intersection of two active faces in a square has two face germs.

\begin{theorem}[Stokes with corners]
\label{p104:M10.06}
On an oriented smooth manifold with corners of positive dimension,
Stokes holds for compactly supported \(C^1\) forms, with the boundary
integral taken once over each codimension-one local face using its
outward-first sign. This includes the usual oriented simplex, boxes
and products such as an interval times a manifold with boundary.
\end{theorem}
\begin{proof}
One intrinsic way to sum the face integrals is to use pairs
\((p,\gamma)\), where \(\gamma\) is a local face germ through \(p\).
A chart supplies one sheet for each of its boundary hypersurfaces,
restricted to that face. The sheets are glued using the induced
transition maps on faces. To see the gluing is locally unambiguous
at a multiple corner, restrict near the given point: on the
relative interior of a local face, depth is one, and the target
face is locally constant there. The relative interior of a
sufficiently small face patch is connected. Its target face is the
one indicated by the derivative at the point, by following a ray
in the relative interior of the source facet and using the
differentiated velocity cone. Continuity then extends the
correspondence across the closure of the patch. Inverse
transitions give inverse correspondences. Charts on the sheets
are orthant charts of dimension \(n-1\). Distinct germs over the
same point have disjoint sheets locally, and germs over distinct
points are separated using the base; a countable base atlas gives
a countable atlas of sheets. Outward-first orientation agrees
across the gluings, because the active normal half-line in the
quotient by the tangent hyperplane is preserved. The projection
to \(M\) is proper: each point has a compact neighborhood in a
coordinate domain, whose inverse image is a finite disjoint
union of its intersections with the closed coordinate faces.
Finitely many such neighborhoods suffice over a compact set.

The local integral on the orthant uses ordinary integration on
its interior, ignoring the coordinate hyperplanes in the bulk.
The compactly supported partition proof applies verbatim: smooth
bumps on Euclidean open sets restrict to the model and the
subordinate gluing construction uses only local compactness.
Use the coefficients in \cref{p104:M10.03}. Full-line derivatives
give zero. A constrained variable \(x_j\) gives the negative
value of \(a_j\) on \(x_j=0\), integrated over the other
variables. In the frame obtained by deleting the \(j\)-th
coordinate, the outward-first sign of that face is \((-1)^j\);
the coefficient there is \((-1)^{j-1}a_j\). Their product gives
the same negative sign. Other summands of the form vanish on
this face. Its intersections with the other faces are lower
dimensional and do not add boundary integrals. Finally localize
by the partition and pull it to the face sheets. The original
compact support lifts to compact support on the sheets, so the
same finite sum gives the stated identity. Equivalently, one
can integrate over the depth-one locus using the restrictions
of these sheet integrals. The support on that open locus alone
need not be compact, even for a compact square.
\end{proof}

For an interval times a boundary manifold the corner charts
are products of the corresponding half-spaces. For the standard
simplex the active affine face functions are independent on their
common intersection, and complete to affine coordinates giving
the corner charts. The pair-of-germs description is used here to
specify the face integral and does not require a more general
choice of category of corner maps. See
\cite[Chapter 16, pp.~415--419]{lee}, the corrections on
pp.~418--419 of \cite{leeerrata}, and \cite{joyce} for other
corner conventions.

\section{Periods and embedded integration}

\begin{corollary}[Periods of a closed one-form]
\label{p104:M10.07}
Integrals of a smooth closed one-form along piecewise smooth paths
are invariant under continuous homotopy with fixed endpoints. On
a simply connected manifold, every closed one-form is exact.
\end{corollary}
\begin{proof}
A piecewise smooth path can be reparametrized on its finitely
many pieces by smooth nondecreasing surjections that are flat at
the endpoints, and can be made constant on endpoint collars.
The joined path is smooth. Substitution in the line integrals
preserves their values. The reparametrized path is homotopic
relative to its endpoints to the original, by interpolating the
parameter functions.

First take the target without boundary. Given a continuous
fixed-endpoint homotopy of the smoothed paths on a square,
precompose in both variables with smooth nondecreasing
surjections of the interval that are constant near both
ends. The resulting map is smooth on a neighborhood in the
square of its four edges, and is constant near the two
side edges at the path endpoints. Extend it continuously to
an open neighborhood of the square by clamping the
parameters to the interval. It is smooth near the boundary
of the square because it was made independent of the
normal parameter there. Relative smooth approximation fixes
a closed collar of the boundary, producing a smooth
homotopy with the same edge data, up to the harmless
reparametrizations. Apply corner Stokes to the pullback
of the closed form. The side integrals are zero, and the
other two terms have opposite signs.

For a target with boundary, use the global smooth homotopy
\(r_t\) of \cref{p104:M06.06}, from the identity to a map
\(r_1\) into the interior. Apply the preceding argument to
the fixed-endpoint homotopy postcomposed with \(r_1\).
For each smoothed path \(\gamma\), apply corner Stokes
to \((s,t)\mapsto r_t(\gamma(s))\), with orientation
\(\dd t\wedge\dd s\). The integral on \(r_1\circ\gamma\)
minus that on \(\gamma\) is the integral on the terminal
endpoint trajectory \(t\mapsto r_t(\gamma(1))\) minus that
on the initial endpoint trajectory. The two paths have
the same endpoints, so these corrections agree. This
proves invariance also in the boundary case. If the manifold is
simply connected, all loop periods vanish. The period
criterion for exactness in \cref{p104:M07.05} then gives
a smooth primitive.
\end{proof}

\begin{proposition}[Compatibility with the embedded D04 model]
\label{p104:M10.08}
Let a smooth embedded submanifold of Euclidean space, possibly with
boundary, be equipped with the smooth structure from its embedded
charts. For an orientation and form in the common domain of the
abstract and D04 constructions, the induced boundary orientation,
restriction and intrinsic exterior derivative agree under the tangent
identification, and the top-form and density integrals have the same
values. In particular the abstract Stokes identity specializes to
the intrinsic D04 formula on this common domain.
\end{proposition}
\begin{proof}
Let \(c:U\to M\) be an embedded parametrization. The identification
in \cref{p104:M02.05} sends the abstract coordinate vector determined
by \(e_j\) to \(Dc(e_j)\). Consequently the \(I\)-th coefficient of
the abstract pullback is the value of the embedded tangent form on
the ordered list \(Dc(e_j)\) for \(j\in I\). In the D04 intrinsic
derivative the chart coefficient is obtained by differentiating
these same coordinate coefficients and antisymmetrizing. The
two derivative formulas therefore agree in the interior and,
by continuity, on the face. The support agrees because the
tangent identification is pointwise an isomorphism and the
base homeomorphism preserves closure.

In a boundary chart the tangent hyperplane of the boundary is the
image of the tangential coordinate directions. The chosen outward
conormal in the embedded definition represents the same outward
half-line modulo this hyperplane. Replacing a representative of
that half-line by another multiplies its normal component by a
positive number and adds a tangent vector. Neither changes the
outward-first orientation. The boundary restriction is evaluation
after the same tangent inclusion. In dimension one this comparison
uses the signed zero-dimensional boundary, with the positive end
of a positively parametrized interval contributing its value
and the negative end subtracting its value.

Take any D04 finite subordinate chart partition on a neighborhood
of the support, with its compatible chart signs. A localized
top coefficient is exactly the coefficient just computed times
the restricted ambient weight. For a D04 density field with
scalar coefficient \(\rho\), the abstract section is
\(\rho\mu_g\), with \(\mu_g\) the induced Euclidean volume
density. The D04 Gram factor evaluates \(\mu_g\) on
\(Dc(e_1),\ldots,Dc(e_n)\). Multiplying by \(\rho(c(u))\)
and the partition weight \(w(c(u))\) gives exactly the
localized coefficient on both sides. A
localized coefficient is continuous on its relative chart and
has compact support away from the artificial edges. Its zero
extension is bounded with compact support and possible
discontinuities only on the coordinate boundary hyperplane.
It is Riemann integrable in a support-enclosing box, with the
same value as its Lebesgue integral, by the Riemann criterion
and agreement of the two integrals there. Full-dimensional
charts without a face give continuous zero extensions.
The finite sums hence coincide, and independence of localization
on each side gives equality of the integrals. For example, on a
smooth oriented embedded surface this identifies the oriented
top-form coefficient and the unsigned Gram density separately;
in dimension zero it identifies the point sums directly.
\end{proof}

The D04 statements used for this comparison are the embedded
chart and tangent interfaces and the intrinsic derivative and
integration interfaces of Chapters 14--15 of \cite{d04}.
The intrinsic D04 Stokes statement admits \(C^2\) embedded charts
with an intrinsic \(C^1\) form. Its separate ambient \(C^1\)
formulation restricts the ambient exterior derivative. The
comparison above concerns the shared smooth embedded setting
and does not infer an intrinsic derivative from a \(C^1\)
chart alone.
\chapter{De Rham cohomology}

\section{The complex and the homotopy operator}

\begin{definition}[The de Rham algebra]
\label{p104:M11.01}
Write \(\Omega^k(M)\) for the smooth real \(k\)-forms and take it to
be zero for \(k<0\). The de Rham group in degree \(k\) is
\[
 H^k_{\mathrm{dR}}(M)
    =\ker(\dd:\Omega^k\to\Omega^{k+1})/
      \operatorname{im}(\dd:\Omega^{k-1}\to\Omega^k).
\]
A closed form is an element of the kernel and an exact form is an
element of the image. Wedge induces a graded algebra, and smooth
maps induce contravariant maps of graded algebras.
\end{definition}

Indeed \(\dd^2=0\). For closed \(\eta\),
\(\dd(\beta\wedge\eta)=\dd\beta\wedge\eta\), so replacing the first factor
of a product by a cohomologous representative changes the product
by an exact form. Graded commutativity gives the same for the
second factor. Naturality of \(\dd\) and of the wedge product then
gives the contravariant action. In degree zero \(\dd f=0\) says
precisely that \(f\) is locally constant, by the fundamental
theorem in small connected coordinate neighborhoods. Thus
\(H^0\) consists of independently chosen constants on the
connected components. It is \(\R\) on a nonempty connected
manifold and zero on the empty manifold. Degrees above the
dimension vanish. See \cite[pp.~441--443]{lee} and
\cite[\S24]{tu}.

\begin{theorem}[Homotopy invariance and the Poincar\'e lemma]
\label{p104:M11.02}
Smooth maps \(f_0,f_1:M\to N\) which are continuously homotopic
induce the same pullback on ordinary de Rham cohomology.
Every closed form of positive degree on a star-shaped open subset
of \(\R^n\) is exact. More generally a nonempty contractible smooth
manifold has zero de Rham cohomology in positive degrees.
\end{theorem}
\begin{proof}
For a smooth homotopy \(F:M\times[0,1]\to N\), decompose the
pullback of a form into
\[
 F^*\eta=\alpha_t+\dd t\wedge\beta_t,\qquad
 \alpha_t=i_t^*F^*\eta,\qquad
 \beta_t=i_t^*\iota_{\partial_t}F^*\eta ,
\]
where \(i_t(p)=(p,t)\). This is the unique decomposition into
a form with no \(\dd t\) component and one with \(\dd t\) first.
Define \(K\eta=\int_0^1\beta_t\,\dd t\) by integrating smooth
coefficients. Differentiation in \(M\) commutes with integration
on the fixed compact interval, locally also at a boundary by
extension. Since
\[
 \dd(\alpha_t+\dd t\wedge\beta_t)
   =\dd_M\alpha_t+\dd t\wedge(\partial_t\alpha_t-\dd_M\beta_t),
\]
the fundamental theorem of calculus gives
\[
       \dd_M K\eta+K(\dd_N\eta)=f_1^*\eta-f_0^*\eta .
\]
The formula works in corner charts if \(M\) has boundary, by
the same coefficient calculation. For continuously homotopic
smooth endpoint maps to a boundaryless target, flatten the
time parameter near both endpoints and extend constantly
outside the interval. Relative Whitney approximation fixes
the two closed time collars in \(M\times\R\), yielding a
smooth homotopy between the given endpoints. For a target
with boundary, postcompose the continuous homotopy with
the interior-valued map \(r_1\) from \cref{p104:M06.06}.
The boundaryless argument equates the pullbacks of the
pushed endpoint maps on the cohomology of the interior.
Apply this equality to the restriction of each class
from the target to its interior. The smooth homotopies \(r_t\circ f_i\) equate
each original pullback with the corresponding pushed one
by the chain homotopy identity above.

On a star-shaped domain the straight-line homotopy joins the
identity to the constant map at the star center. In positive
degree the constant pullback is zero, so the chain homotopy
constructs a primitive of any closed form. The same cohomology
argument works for a continuous contraction with smooth
constant and identity maps. Constants in degree zero remain
as described above.
\end{proof}
The operator and its endpoint identity follow
\cite[pp.~443--448]{lee} and \cite[\S\S27,29]{tu}; the restriction
to positive degree in the exactness statement is recorded in
\cite[p.~447]{leeerrata}.

\section{Mayer--Vietoris and computations}

\begin{theorem}[Open-cover Mayer--Vietoris]
\label{p104:M11.03}
For \(M=U\cup V\) with \(U,V\) open, restrictions give a short exact
sequence of de Rham complexes
\[
 0\longrightarrow\Omega^\bullet(M)\xrightarrow{r}
 \Omega^\bullet(U)\oplus\Omega^\bullet(V)\xrightarrow{q}
 \Omega^\bullet(U\cap V)\longrightarrow0,\qquad
 q(a,b)=b|_{U\cap V}-a|_{U\cap V}.
\]
It gives the natural long exact sequence on cohomology, with
connecting map increasing degree by one.
\end{theorem}
\begin{proof}
Smooth forms glue uniquely when they agree on an overlap. To
lift a form \(\gamma\) on \(U\cap V\), take a partition
\(\rho_U+\rho_V=1\) subordinate to \(U,V\). Use
\(a=-\rho_V\gamma\) on \(U\cap V\) and extend by zero over \(U\setminus V\).
Use \(b=\rho_U\gamma\) and extend over \(V\setminus U\). At every
point of a zero-extension locus there is a neighborhood on which
the multiplying factor vanishes, by its closed support
subordination. Thus the extensions are smooth, and \(q(a,b)=\gamma\).

For closed \(\gamma\), \(\dd(a,b)\) has zero difference and is the
restriction of a unique closed form \(\theta\) on \(M\).
Define \(\delta[\gamma]=[\theta]\). Changing the lift by \(r(\nu)\)
changes \(\theta\) by \(\dd\nu\), and a change of representative by
\(\dd\gamma_0\) can be lifted by adding the derivative of a lift of
\(\gamma_0\). To check exactness at the overlap group, a closed
lift gives zero connecting class, and if \(\theta=\dd\nu\), then
subtracting \(r(\nu)\) from the lift makes it closed. At the group
of \(M\), a closed \(\theta\) restricting to an exact pair is the
derivative of a pair \((a,b)\), in the sense \(r(\theta)=\dd(a,b)\).
The difference \(q(a,b)\) is closed and its connecting class is
\([\theta]\). Finally, a closed pair whose difference is exact
can be changed by an exact pair lifting that primitive, and
then glues to a closed form on \(M\). This proves exactness
at every position. Pullbacks commute with the difference, restriction
and differential maps for maps respecting the cover. Applying
them to any chosen lift proves naturality of \(\delta\) without
requiring a natural choice of partition.
\end{proof}
See \cite[pp.~448--451, 460--464]{lee} and \cite[\S\S25--26]{tu}.
This construction will also identify the connecting map in a
surface calculation below.

\begin{proposition}[Spheres, tori, and an oriented surface]
\label{p104:M11.04}
For \(n\geq1\), the only nonzero real de Rham groups of \(S^n\)
are one-dimensional in degrees zero and \(n\). The algebra of
\((S^1)^r\) is the exterior algebra on the \(r\) classes of the
coordinate circle one-forms. A genus-two oriented surface
constructed as a connected sum of two tori has dimensions
\(1,4,1\) in degrees \(0,1,2\).
\end{proposition}
\begin{proof}
For \(S^n\) cover by the complements of the two poles. Both
sets are diffeomorphic to \(\R^n\); their intersection is
\(S^{n-1}\times\R\), which retracts smoothly onto \(S^{n-1}\).
When \(n=1\), the intersection has two components. The
degree-zero difference map has image the diagonal in
\(\R^2\), so its cokernel is \(\R=H^1(S^1)\).
When \(n>1\) the intersection is connected, the degree-zero
difference is surjective, and the long exact sequence
identifies \(H^k(S^n)\) with \(H^{k-1}(S^{n-1})\) for \(k\geq2\).
Induction gives the claim. The description \(H^0(S^0)=\R^2\)
accounts for the initial two components.

To keep track of the product ring, write \(S^1=\R/\mathbb Z\)
with the one-form \(\vartheta\) descended from \(\dd t\).
For any \(X\), circle translation \(R_s\) on \(X\times S^1\)
is smoothly homotopic to the identity through \(R_{us}\),
\(0\leq u\leq1\). The operator in \cref{p104:M11.02} applied
to this homotopy depends smoothly on \(s\). Integrate its
identity over \(0\leq s\leq1\). The averaging projection
\(A=\int_0^1 R_s^*\,\dd s\) is thereby chain-homotopic to the
identity. Its image consists exactly of translation-invariant
forms: translation preserves the full-period integral, and
an invariant form is fixed by the average. Inclusion of
invariant forms induces an isomorphism, since \(A\) is a
chain projection with this homotopy. An invariant form
decomposes uniquely as
\[
       \pi_X^*\alpha+\vartheta\wedge\pi_X^*\beta .
\]
To obtain the coefficients, restrict the form and its insertion
by the unit translation field to \(X\times\{0\}\), then use
invariance. The differential is \(\dd\alpha-\vartheta\wedge \dd\beta\)
with pullbacks understood. This subcomplex is also a subalgebra,
so its cohomology gives the graded ring \(H^*(X)\otimes
\Lambda([\vartheta])\), with the usual sign when the degree-one
generator is commuted across a class. Induction from a point
gives the torus algebra.

For the surface calculation, let \(D_1\subset D_2\) be concentric
coordinate disks in an oriented torus, with the closure of the
smaller contained in the larger. Cover the torus by
\(U=T^2\setminus\overline{D_1}\) and \(V=D_2\).
The overlap is an annulus, with generator \(\gamma\) of \(H^1\)
given by the angular one-form of period one. Choose \(\rho_V\)
radial, equal to one near the inner edge of the annulus and
zero near its outer edge, and use the lift in the preceding
proof. The connecting two-form on the overlap is
\(-\dd\rho_V\wedge\gamma\) and is zero off the transition
annulus. Its integral is one in the corresponding orientation,
by integrating in the radial and angular variables. In
particular it is not exact by Stokes on the closed torus.
Since \(H^2(T^2)\) is one-dimensional, the annulus connecting
map is an isomorphism. Exactness now shows that
\(H^1(T^2)\to H^1(U)\) is an isomorphism, that the restriction
\(H^1(U)\to H^1(U\cap V)\) is zero, and that \(H^2(U)=0\).

Glue two tori with open disks removed along their boundary
circles by a reflection in angular coordinates. Use signed
radial collars to chart the seam by a product of the circle
with an open interval. Together with the interior charts
these define a smooth surface: transitions away from the
seam are collar transitions, and the reflection reverses
boundary orientation so that the two bulk orientations
agree in the signed collar. The quotient is Hausdorff,
since its equivalence relation on the compact disjoint
union is closed; the two interiors and the collar supply
a countable chart cover.

Cover this connected sum by the two sides each extended
slightly across the seam. Their intersection is an open
annulus. Each side is diffeomorphic to \(U\) above.
In the collar it suffices to identify a slightly extended
radial interval \((-\epsilon,a)\) with \((0,a)\), by a
smooth increasing diffeomorphism that is the identity
near \(a\). For example, integrate
\(1-\epsilon\chi/\int\chi\) from \(-\epsilon\), with
a smooth \(0\leq\chi\leq1\) supported in the lower part
of this interval and \(\int\chi>\epsilon\).
The function has positive derivative and equals the
identity above the support. With small collar sizes,
the intersection maps to a concentric subannulus of the
original punctured torus collar. The annulus restriction
in degree one is zero on both sides by the preceding
calculation, and degree two vanishes on each side.
The degree-zero difference is surjective since the
intersection is connected. Mayer--Vietoris gives the
direct sum of the two degree-one groups and one copy of
\(\R\) in degree two.
\end{proof}
These are applications of the sphere and product computations
in \cite[Chapter 17]{lee} and \cite[\S28]{tu}. The surface
calculation uses the degree-one connecting map just constructed
and the already proved Stokes theorem.

\section{Compact supports and top-degree integration}

\begin{proposition}[The compact-support complex]
\label{p104:M11.05}
The spaces \(\Omega_c^\bullet(M)\) of compactly supported smooth
forms form a subcomplex. A smooth proper map induces pullback
on its cohomology. An open cover \(M=U\cup V\) gives the
compact-support exact sequence
\[
 0\longrightarrow\Omega_c^\bullet(U\cap V)
 \xrightarrow{(\operatorname{ext}_0,-\operatorname{ext}_0)}
 \Omega_c^\bullet(U)\oplus\Omega_c^\bullet(V)
 \xrightarrow{\operatorname{ext}_0+\operatorname{ext}_0}
 \Omega_c^\bullet(M)\longrightarrow0 .
\]
In particular it gives the long exact sequence in compact-support
cohomology.
\end{proposition}
\begin{proof}
The derivative is local, hence does not enlarge support.
The support of \(f^*\eta\) is contained in \(f^{-1}(\supp\eta)\),
compact when \(f\) is proper. Extension by zero of a form
supported in a compact subset of an open submanifold is smooth
and commutes with \(\dd\). A pair in the kernel of the sum has
equal opposite extensions. Their common support in \(M\) is
compact and lies in \(U\cap V\), so the pair comes from the
first map. For surjectivity split the compactly supported
form by a subordinate partition. Its two pieces have compact
support in their respective open sets. The long exact
sequence follows by the same closed-lift argument as above.
\end{proof}
Unlike the ordinary sequence, the compact-support sequence
uses extension by zero in both of its displayed nontrivial
arrows. See \cite[pp.~451--455]{lee}.

\begin{theorem}[Compact-support Poincar\'e calculation]
\label{p104:M11.06}
For \(n\geq0\), \(H_c^k(\R^n)\) is zero for \(k\neq n\) and is
isomorphic to \(\R\) in degree \(n\) by integration. In
particular a compactly supported top form has a compactly
supported primitive if and only if its total integral is zero,
with exactness in dimension zero understood in the complex.
\end{theorem}
\begin{proof}
For \(\R\times X\) write \(\eta=\alpha_t+\dd t\wedge\beta_t\)
and put \(I\eta=\int_{\R}\beta_t\,\dd t\). Fix a smooth
compactly supported \(\rho\) on \(\R\) of integral one and put
\(J\gamma=\rho(t)\,\dd t\wedge\gamma\). We have
\(I J=\id\), \(I\dd=-\dd I\), and \(\dd J=-J\dd\); for \(I\dd\), the
integral of \(\partial_t\alpha_t\) vanishes by compact support.
Set
\[
 B\eta(t)=\int_{-\infty}^t
           \bigl(\beta_s-\rho(s)I\eta\bigr)\,\dd s .
\]
The integral is zero for \(t\) below or above a sufficiently
large fixed compact interval, the latter because the corrected
integrand has total integral zero. The support in \(X\) lies
in the compact projection of \(\supp\eta\). Thus \(B\eta\)
is compactly supported and smooth. Directly,
\[
  B(\dd\eta)=\alpha_t-\dd_XB\eta,\qquad
  \dd(B\eta)=\dd_XB\eta+\dd t\wedge(\beta_t-\rho(t)I\eta),
  \qquad \dd B+B\dd=\id-JI.
\]
The first identity uses \(I(\dd\eta)=-\dd_XI\eta\). These formulas
include degrees zero and one with the absent negative-degree
components interpreted as zero. Hence \(I\) and \(J\) give
inverse cohomology isomorphisms with a shift of one degree.
Apply this successively to the Euclidean coordinates starting
with a point, whose only group is \(H_c^0=\R\).
With the new coordinate placed first in the orientation at
each step, the resulting scalar is the usual total integral
by Fubini.
\end{proof}
The support correction is essential. The positive-degree
ordinary radial homotopy does not have this support property.
See \cite[Lemma 17.27 and Theorem 17.28]{lee}.

\begin{theorem}[Top-degree cohomology]
\label{p104:M11.07}
Let \(M\) be a nonempty connected \(n\)-manifold without boundary.
If it is oriented, integration gives an isomorphism
\(H_c^n(M)\to\R\). Thus \(H^n(M)=\R\) if \(M\) is compact.
If \(M\) is noncompact, \(H^n(M)=0\), independently of
orientability. For nonorientable \(M\), \(H_c^n(M)=0\) as well.
\end{theorem}
\begin{proof}
For the oriented compact-support assertion, integration vanishes
on exact compactly supported forms by Stokes. An oriented coordinate
ball contains a bump top form of integral one. For a compactly
supported top form \(\eta\), use a partition into finitely many
terms each compactly supported in an oriented coordinate ball.
In each ball, subtract the term's integral times a chosen
unit-integral form. The difference has compactly supported
primitive by \cref{p104:M11.06}, since a coordinate ball is
orientation-preservingly diffeomorphic to \(\R^n\). Extend the
primitive by zero. To compare two unit-integral forms on balls
with nonempty intersection, choose a third such form in a
smaller ball in the intersection. Subtract it from each of the
two forms in turn and use the same zero-integral argument.
The intersection graph of a cover by connected coordinate
balls is connected because a disconnected graph would
separate \(M\) into disjoint open unions. Any two balls
can thus be connected by a finite chain. All unit-integral
classes agree, and the original form represents its integral
times that class. In dimension zero the oriented connected
space is a point, with signed evaluation, and the statement
is immediate.

Here is the exhaustion step for the ordinary noncompact
assertion in the oriented case. Choose a countable locally
finite open cover by connected precompact coordinate balls.
One construction is to use a compact exhaustion
\(K_j\subset\operatorname{int}K_{j+1}\) and cover the compact
shells \(K_j\setminus\operatorname{int}K_{j-1}\) by finitely
many balls with closure in
\(\operatorname{int}K_{j+1}\setminus K_{j-2}\), omitting
negative indices. Its intersection graph is connected,
infinite and locally finite, the latter since a compact
closure meets only finitely many cover elements. Enumerate
its vertices. At each step, take the first remaining vertex
and group it with all finite connected components that become
detached in the remaining graph after deleting it. Every
component of the graph just before that step is infinite.
A newly finite component meets a neighbor of the deleted
vertex, so only finitely many such components are added.
The resulting group is finite and connected, and the graph
of remaining vertices again has only infinite components.
Every vertex is eventually grouped by the choice of the
first remaining one. The unions \(V_j\) of the balls in each
group are connected precompact open sets. They form a locally
finite cover, since the original cover is locally finite
and each original ball was assigned to one group. Each group
has an edge to a later group, since it is finite inside an
infinite component at its removal step. Thus there is
\(k(j)>j\) with \(V_j\cap V_{k(j)}\neq\varnothing\).

Choose a partition subordinate to the \(V_j\) and write an
arbitrary top form as \(\eta=\sum_j\eta_j\), each piece
compactly supported in \(V_j\). Choose an \(n\)-form
\(\theta_j\) of integral one compactly supported in a
coordinate ball in \(V_j\cap V_{k(j)}\). Inductively set
\[
     c_j=\int_{V_j}\left(\eta_j+
              \sum_{i:\,k(i)=j} c_i\theta_i\right).
\]
Only \(i<j\) occur in this sum. The compact-support
assertion on the connected oriented open set \(V_j\) gives
\(\beta_j\in\Omega_c^{n-1}(V_j)\) with
\[
   \dd\beta_j=\eta_j+\sum_{i:\,k(i)=j}c_i\theta_i-c_j\theta_j.
\]
The extensions by zero of the \(\beta_j\) form a locally
finite sum on \(M\). In the derivative of that sum each
transfer term cancels with its contribution in the unique
later group. Thus \(\dd(\sum_j\beta_j)=\eta\). This explicitly
uses connected precompact sets at every application of
the compact-support assertion.

For nonorientable \(M\), use its connected orientation double
cover \(\pi:\widetilde M\to M\) with orientation-reversing
deck involution \(\tau\). The finite covering is proper:
a compact set has a finite cover by precompact evenly
covered neighborhoods with closures inside larger evenly
covered sets, giving finitely many compact lifted
neighborhood closures. If \(\eta\) has compact support,
the compactly supported form \(\pi^*\eta\) is invariant
under \(\tau^*\), so its oriented integral equals its
negative and vanishes. Take its compactly supported
primitive upstairs. Averaging it with its pullback by
\(\tau\) yields an invariant primitive. Such a form
descends in local sheets, where invariance ensures the
two definitions coincide. Its projected support is
compact. If \(M\) is noncompact, its cover is noncompact
as well; the ordinary oriented result yields a primitive
without support restrictions, and the same averaging
gives descent. In the compact case ordinary and
compact-support complexes coincide. This proves all
the assertions.
\end{proof}
For the connected exhaustion argument, compare the corrected
proof of Theorem 17.32 at pp.~455--456 of \cite{leeerrata}.
The compact-support argument and orientation-cover comparison
correspond to \cite[Theorems 17.30--17.34]{lee}.
\chapter{Degree and the comparison map}

\section{Degree of a smooth map}

Here a closed manifold is compact and without boundary. For a
nonempty closed connected oriented \(n\)-manifold, top-degree
cohomology has the distinguished generator of integral one. A
smooth map between two such manifolds acts on this generator
by a real scalar. The following proof identifies the scalar
as an integer.

\begin{theorem}[Degree and signed regular fibers]
\label{p104:M12.01}
Let \(F:M\to N\) be smooth between nonempty closed connected
oriented \(n\)-manifolds. There is an integer \(\deg(F)\) such that
\[
                 \int_M F^*\eta=\deg(F)\int_N\eta
                   \qquad(\eta\in\Omega^n(N)).
\]
For every regular value \(y\in N\) the fiber is finite, and
\[
       \deg(F)=\sum_{x\in F^{-1}(y)}
                     \operatorname{sign}(D F_x).
\]
Here the sign is that of the tangent isomorphism between the
specified orientations. An empty fiber contributes zero.
\end{theorem}
\begin{proof}
The integral of the pullback vanishes on exact forms by
naturality of \(\dd\) and Stokes. By \cref{p104:M11.07}, it
therefore equals a scalar times the target integral.
For positive dimension the inverse function theorem
makes the fiber at a regular value discrete. The fiber is
closed in compact \(M\) and is consequently finite. Choose
pairwise disjoint neighborhoods \(U_x\) of the fiber points
on which \(F\) is a diffeomorphism onto neighborhoods of \(y\).
The compact complement of their union has image missing
\(y\). Choose a small connected oriented coordinate ball
\(V\) about \(y\) contained in all the image neighborhoods
and disjoint from the image of that complement. For an empty
fiber take \(V\) disjoint from \(F(M)\). Choose \(\eta\) supported
compactly in \(V\) with integral one. The preimage of \(V\)
is the disjoint union of the corresponding inverse
neighborhoods in the \(U_x\). The pullback is supported
compactly in their union, and by signed substitution its
integral on each inverse neighborhood is the orientation
sign of \(D F_x\). This proves the formula. Sard's theorem
ensures regular values exist: a nonempty target chart of
positive dimension cannot be covered by a null set. For
dimension zero each connected manifold is a point and
its integral is evaluation times its orientation sign,
giving the stated scalar and sign directly.
\end{proof}
The degree here is an oriented count, so it is independent of
the chosen regular value even when different regular fibers
have different cardinalities. A nonzero degree forces
surjectivity, since a point outside the image is a regular value
with empty fiber. See \cite[Theorems 17.35--17.36]{lee}.

\begin{corollary}[Homotopy, composition and boundary extension]
\label{p104:M12.02}
Under the hypotheses of \cref{p104:M12.01}, degree is invariant
under continuous homotopy through maps with smooth endpoints,
the identity has degree one, and degrees multiply under
composition. If \(W\) is a compact oriented smooth
\((n+1)\)-manifold with boundary and a smooth \(G:W\to N\)
extends a map on its boundary, the sum of the degrees of
the restrictions to the boundary components with their
induced orientations is zero.
\end{corollary}
\begin{proof}
On top cohomology pullbacks are functorial and homotopy
invariant by \cref{p104:M11.02}. The integral generators turn
these identities into the assertions about degree. For
the boundary statement choose a smooth \(n\)-form
\(\eta\) of integral one on \(N\). Since \(\dd\eta=0\),
Stokes gives
\[
 \int_{\partial W}(G|_{\partial W})^*\eta
       =\int_W \dd(G^*\eta)=\int_W G^*(\dd\eta)=0 .
\]
The boundary components are open in the boundary and there
are finitely many by compactness. The left side is their
signed degree sum. The empty boundary gives the empty sum.
\end{proof}
For example, the identity of \(S^{n}\), \(n\geq1\), cannot
extend smoothly across the ball with the induced sphere
orientation. It also cannot extend continuously. Indeed,
if a continuous extension existed, its values near the
sphere would be uniformly close to the radial projection,
since the extension agrees with the identity on the sphere
and the ball is compact. Interpolate there with the radial
projection in the ambient Euclidean space and normalize to
unit length, using a smooth radial cutoff equal to one
close to the sphere and zero away from that neighborhood.
Uniform closeness ensures the vector never vanishes.
The resulting continuous map is smooth on a collar of the
sphere. Relative Whitney approximation fixes a smaller
closed collar and gives a smooth extension, which was
excluded. For \(S^0\) the same nonextension follows from
connectedness of the interval. This is the usual no-retraction
input to Brouwer's theorem; the equivalent no-retraction
and fixed-point statements also appear in \cite[Chapter 17]{d10}.
See \cite[pp.~457--460]{lee} for the degree route.

\section{Smooth simplices and the integration cochain map}

The standard \(q\)-simplex is
\[
 \Delta^q=\{(t_1,\ldots,t_q):t_i\geq0,\ 
                         t_1+\cdots+t_q\leq1\},
\]
with ordered vertices \(v_0=0,v_i=e_i\) and the orientation of
\(\dd t_1\wedge\cdots\wedge\dd t_q\). At a vertex or a higher
face intersection it has the corner structure of
\cref{p104:M10.06}. A smooth singular simplex in \(M\) is a
map \(\sigma:\Delta^q\to M\) smooth up to the faces, using
local Euclidean extensions in charts. Let
\(C_q^{\mathrm{sm}}(M;\R)\) be the real vector space freely
generated by those simplices, with boundary
\[
 \partial\sigma=\sum_{j=0}^q(-1)^j\sigma\circ\varepsilon_j,
 \qquad
 \varepsilon_j:\Delta^{q-1}\longrightarrow\Delta^q
\]
the affine parametrization of the face obtained by omitting
vertex \(j\) and preserving the remaining order. Degree zero
has boundary zero. The two ways of omitting a fixed pair
of vertices yield the same affine map with opposite signs,
so \(\partial^2=0\). The dual differential on real cochains
is \((\delta c)(\sigma)=c(\partial\sigma)\).

\begin{proposition}[The natural integration cochain map]
\label{p104:M12.03}
The rule
\[
     \mathcal I_M(\omega)(\sigma)=\int_{\Delta^q}\sigma^*\omega
          \qquad(\omega\in\Omega^q(M))
\]
defines a natural cochain map from the de Rham complex to
the real smooth singular cochain complex. It induces a
natural map \(H_{\mathrm{dR}}^q(M)\to H_{\mathrm{sm}}^q(M;\R)\).
\end{proposition}
\begin{proof}
The simplex is compact, so the pulled-back form is
integrable with the conventions of the integration chapter.
To verify its face signs, for \(j\geq1\) take outward
vector \(-e_j\) along \(t_j=0\). The positively ordered
face chart is based at \(v_0\) and has the coordinate
basis omitting \(e_j\). Placing \(-e_j\) first gives
determinant sign \((-1)^j\). On the face opposite \(v_0\)
one can use the outward vector \(e_1\) and face basis
\(e_i-e_1\), \(2\leq i\leq q\), which together have
positive determinant. The formulas include \(q=1\)
as the signed endpoint rule. Therefore
\[
 \mathcal I_M(\dd\omega)(\sigma)
  =\int_{\Delta^{q+1}}\dd(\sigma^*\omega)
  =\sum_j(-1)^j\int_{\Delta^q}(\sigma\circ\varepsilon_j)^*\omega
  =(\delta\mathcal I_M(\omega))(\sigma)
\]
for a \((q+1)\)-simplex, by naturality of the derivative
and corner Stokes. Finally for smooth \(F:M\to N\),
evaluating \(\mathcal I_M(F^*\omega)\) on \(\sigma\)
is the same as evaluating \(\mathcal I_N(\omega)\) on
\(F\circ\sigma\). These identities prove the claim.
\end{proof}

\paragraph{Comparison theorem (statement).}
Let \(C_\bullet(M;\R)\) use all continuous singular
simplices. Restriction of cochains from continuous
to smooth singular simplices induces an isomorphism
\(r^*:H^q_{\mathrm{sing}}(M;\R)\to H^q_{\mathrm{sm}}(M;\R)\).
The de Rham theorem asserts that
\((r^*)^{-1}\mathcal I_M^*\) is a natural isomorphism of
real vector spaces \(H_{\mathrm{dR}}^q(M)\cong
H^q_{\mathrm{sing}}(M;\R)\) for each \(q\).
The two assertions beyond \cref{p104:M12.03} require
the comparison of smooth and continuous chains and
the global comparison argument. Their full proofs
belong to the algebraic-topology comparison unit.
A proof needs a subdivision or small-chains theorem,
relative smoothing compatible with faces, and a
local-to-global argument for the comparison on
cohomology. In particular neither approximation of
isolated simplices nor the Stokes calculation alone
gives either isomorphism. The construction and comparison
are located in \cite[Theorem 18.12, Proposition 18.13,
and Theorem 18.14]{lee}.

\backmatter
\begin{thebibliography}{99}
\phantomsection
\addcontentsline{toc}{chapter}{Bibliography}
\bibitem{lee} John M. Lee, \href{https://doi.org/10.1007/978-1-4419-9982-5}{\emph{Introduction to Smooth Manifolds}}, 2nd ed., Graduate Texts in Mathematics 218, Springer, 2013.
\bibitem{tu} Loring W. Tu, \href{https://doi.org/10.1007/978-1-4419-7400-6}{\emph{An Introduction to Manifolds}}, 2nd ed., Universitext, Springer, 2011.
\bibitem{leeerrata} John M. Lee, \href{https://sites.math.washington.edu/~lee/Books/ISM/errata.pdf}{Corrections to \emph{Introduction to Smooth Manifolds}, 2nd ed.}
\bibitem{tuerrata} Loring W. Tu, \href{https://lwtu.github.io/doc/mf2_master_errata2.pdf}{Corrections to \emph{An Introduction to Manifolds}, 2nd ed.}
\bibitem{joyce} Dominic Joyce, \href{https://arxiv.org/abs/0910.3518}{\emph{On manifolds with corners}}, arXiv:0910.3518.
\bibitem{d04} MathGenome, \emph{Unified Multivariable Analysis}, D04, in MathGenome v0.8.0.
\bibitem{d05} MathGenome, \emph{Unified Linear Algebra}, D05, in MathGenome v0.8.0.
\bibitem{d10} MathGenome, \emph{Unified Point-Set Topology}, D10, in MathGenome v0.8.0.
\bibitem{d11} MathGenome, \emph{Unified Real Analysis and Measure Integration}, D11, in MathGenome v0.8.0.
\bibitem{d13} MathGenome, \emph{Unified Ordinary Differential Equations and Dynamical Systems}, D13, in MathGenome v0.8.0.
\end{thebibliography}
\end{document}
